% Compile from this directory with:
%   latexmk -pdf -interaction=nonstopmode final_connected_pseudospectrum_proof_2026-07-16.tex
% Regenerate all figures with:
%   python3 make_hatano_nelson_figures.py
\documentclass[review,hidelinks,onefignum,onetabnum]{siamart251216}

\usepackage[T1]{fontenc}
\usepackage{amsfonts,amssymb,mathtools}
\usepackage{booktabs}
\usepackage{graphicx}
\usepackage{microtype}

\allowdisplaybreaks
\graphicspath{{figures/}}
\DeclareGraphicsExtensions{.pdf,.png,.jpg}

\newsiamremark{remark}{Remark}
\DeclareMathOperator{\tr}{tr}
\DeclareMathOperator{\diag}{diag}
\DeclareMathOperator{\tridiag}{tridiag}

\newcommand{\C}{\mathbb C}
\newcommand{\R}{\mathbb R}
\newcommand{\ii}{\mathrm i}

\headers{Connected Pseudospectra: Open Hatano--Nelson Chain}{S. Jin}
\title{Connected Pseudospectra of a Nonnormal Toeplitz Path:
The Open Hatano--Nelson Chain\thanks{{\bfseries Funding:} None.
{\bfseries Conflicts of Interest:} None.
{\bfseries Prior Publication:} None.}}
\author{Shanmu Jin\thanks{Department of Neurosurgery, Peking Union Medical
College Hospital, No.~1 Shuaifuyuan, Dongcheng District, Beijing 100730,
China (\email{jinshanmu1129@pku.edu.cn}).}}

\ifpdf
\hypersetup{
  pdftitle={Connected Pseudospectra of a Nonnormal Toeplitz Path: The Open Hatano--Nelson Chain},
  pdfauthor={Shanmu Jin}
}
\fi

\begin{document}
\maketitle

\begin{abstract}
For the nonnormal tridiagonal Toeplitz matrix
$A_n(a)=\tridiag(a,0,1)$, $0<a<1$, we determine the finite-size
transition at which its Euclidean $\varepsilon$-pseudospectrum becomes
connected.  The least singular value of $(x+\ii y)I_n-A_n(a)$ is
nondecreasing in $|y|$, so every pseudospectral component retracts onto
an interval of the real axis and is therefore contractible.  If
$\gamma_n$ is the largest least-singular-value barrier between adjacent
open-boundary eigenvalues, then connectedness is equivalent to
$\varepsilon>\gamma_n$.  We prove the strict finite-size comparison
$\gamma_{n+1}<\gamma_n$ and sharp-order bounds
$\gamma_n\asymp_a(n+1)a^{n/2}$.  Consequently, the first connected
matrix size is also the eventual connectedness threshold, and, for fixed
$a$, it equals $2(\log(1/\varepsilon)+\log\!\log(1/\varepsilon))/
\log(1/a)+O_a(1)$ as $\varepsilon$ tends to zero.  In the open-boundary
Hatano--Nelson model, the coefficient $2/\log(1/a)$ is the amplitude skin
depth of the right eigenmodes, in lattice-spacing units.  The threshold is
the site count at which fixed unstructured uncertainty exceeds every
inter-eigenvalue barrier: below it, disconnected components retain
perturbation-stable spectral labels, whereas at and above it they form one
worst-case uncertainty cloud.
\end{abstract}

\begin{keywords}
pseudospectrum, nonnormal Toeplitz matrix, Hatano--Nelson model,
non-Hermitian skin effect, least singular value, Chebyshev polynomial
\end{keywords}

\begin{MSCcodes}
15A18, 15A60, 47A10, 65F15, 81Q12
\end{MSCcodes}

\section{Introduction}\label{sec:introduction}

For a matrix $A\in\C^{n\times n}$, the Euclidean
$\varepsilon$-pseudospectrum is
\begin{equation}\label{eq:intro-pseudospectrum}
 \sigma_\varepsilon(A)
 =\sigma(A)\cup
 \{z\in\C:\lVert(zI_n-A)^{-1}\rVert_2>\varepsilon^{-1}\}
 =\{z\in\C:s_{\min}(zI_n-A)<\varepsilon\}.
\end{equation}
Equivalently, it is the union of the spectra of $A+\Delta A$ over all
unstructured complex perturbations $\Delta A\in\C^{n\times n}$ with
$\lVert\Delta A\rVert_2<\varepsilon$; see
\cite{TrefethenEmbree2005}.  For nonnormal
matrices, these sets reveal spectral sensitivity and resolvent growth
that the eigenvalues alone can miss; see
\cite{Trefethen1997,TrefethenEmbree2005}.  Classical spectral limit sets
for banded Toeplitz sections go back to Schmidt and Spitzer
\cite{SchmidtSpitzer1960}.  Finite Toeplitz matrices are a setting in
which the distinction between spectra and pseudospectra can be
exponentially large in the dimension; their finite-section resolvents
and pseudospectra have been studied in
\cite{BottcherEmbreeTrefethen2002,BottcherGrudsky2005,
BottcherGrudskySilbermann1997,ReichelTrefethen1992}.

We study the real tridiagonal Toeplitz path
\begin{equation}\label{eq:A-matrix}
 A_n(a)=\tridiag(a,0,1)
 =\begin{bmatrix}
 0&1\\
 a&0&\ddots\\
 &\ddots&\ddots&1\\
 &&a&0
 \end{bmatrix},
 \qquad n\geq2,\quad 0<a<1,
\end{equation}
where the three arguments of $\tridiag$ denote the subdiagonal,
diagonal, and superdiagonal, respectively.  Closed-form spectral,
conditioning, and pseudospectral results for general tridiagonal
Toeplitz matrices are developed in
\cite{NoschesePasquiniReichel2013}.  Although $A_n(a)$ is
diagonally similar to a symmetric Jacobi matrix, the similarity is not
unitary, and hence it does not preserve the Euclidean pseudospectrum.
Our question is finite-dimensional and topological: as $n$ grows at
fixed $a$ and $\varepsilon$, can connectedness appear and then disappear, or is
there a unique transition?

The same family is the clean open-boundary Hatano--Nelson chain
\cite{HatanoNelson1996,HatanoNelson1997}.  With the site convention in
which $t_R$ moves a particle from $j$ to $j+1$ and $t_L$ moves it from
$j+1$ to $j$, assume $0<t_R<t_L$ and use the positive-sign one-particle
representative
\[
 H_n^{\rm OBC}=t_LA_n(a),\qquad a=\frac{t_R}{t_L}.
\]
The dimensionless logarithmic parametrization used in the proof is
\[
 h=-\frac12\log a>0,\qquad a=e^{-2h},\qquad
 r=\sqrt a=e^{-h}.
\]
The diagonal gauge
\[
 \mathcal G_n(\varrho)=\diag(1,\varrho,\ldots,\varrho^{n-1})
\]
and the upper-shift matrix \(V_n=\tridiag(0,0,1)\) give the
dimensionless similarity
\[
 A_n(a)=\mathcal G_n(r)\bigl[r(V_n+V_n^T)\bigr]
 \mathcal G_n(r)^{-1}.
\]
Multiplying by \(t_L\) and using \(t_Lr=\sqrt{t_Lt_R}\) recovers the
usual physical imaginary-gauge formula for \(H_n^{\rm OBC}\).

The endpoint \(a=0\) is excluded.  There \(A_n(0)=V_n\) is the
unidirectional Jordan-block limit, the diagonal gauge is singular, and the
limits \(a\downarrow0\) and \(\delta\downarrow0\), where \(\delta\) is the
normalized uncertainty level used below, need not commute.  In particular,
the fixed-\(a\) small-\(\delta\) asymptotic is not extrapolated through
\(a=0\).

The open-boundary eigenvalues are real, while the right eigenmodes
have the envelope $e^{-h(j-1)}$ and accumulate at the left edge.  This
is the non-Hermitian skin effect
\cite{OkumaEtAl2020,YaoWang2018,YokomizoMurakami2019}; see
\cite{ZhangEtAl2022} for a review.  Under periodic boundary conditions
the Bloch dispersion traces the continuous ellipse
\[
 \mathcal E_{\rm PBC}(k)=t_Le^{\ii k}+t_Re^{-\ii k}
 =(t_L+t_R)\cos k+\ii(t_L-t_R)\sin k.
\]
A finite periodic ring samples this curve at \(k=2\pi q/n\),
\(q=0,\ldots,n-1\), whereas the displayed ellipse is obtained as \(k\)
traverses the Brillouin zone.

Early work connecting Hatano--Nelson--type localization with
pseudospectra treated random bidiagonal and Toeplitz models
\cite{TrefethenContediniEmbree2001}.  Direct pseudospectral
interpretations of nonnormal boundary sensitivity appear in
\cite{OkumaSato2020,OkumaSato2021}.  Mathematical foundations for the
skin effect were developed in \cite{AmmariEtAl2024}, while spectra and
pseudospectra of tridiagonal $k$-Toeplitz matrices were treated in the
original article \cite{AmmariEtAl2025}; the affected finite-section
pseudospectral statements were corrected in the corrigendum
\cite{AmmariCorrigendum2026}.

For asymmetric lattices,
Kiorpelidis and Makris \cite{KiorpelidisMakris2025} observed that
eigenvalue-centered pseudospectral clouds expand with the lattice size and
merge, beyond a critical size denoted by \(\widetilde N\), into a single
cloud.  In their description the pseudospectral abscissa grows
exponentially up to that crossover, whereas long-time dynamical
amplification need not stop there.  We write \(N_c\) informally for the first
site count at which the clean-chain Euclidean pseudospectrum is connected;
\cref{thm:main} gives the precise definition.  For this clean chain, \(N_c\)
is a rigorous finite-dimensional analogue of their cloud-merger scale: we
identify the exact finite-size criterion, rule out re-entrant splitting, and
determine the asymptotic scale.  Sirker \cite{Sirker2026} analyzes spectral
instability and Toeplitz-index structure through exponentially small singular
values of \(H_n^{\rm OBC}-EI_n\) at fixed values of the spectral parameter
\(E\) for the same open Hatano--Nelson chain.  The present work instead treats
the global finite-dimensional barrier obtained by optimizing over all
inter-eigenvalue gaps; it determines the topology of the full Euclidean
pseudospectrum and the barrier's strict decrease from dimension \(n\) to
\(n+1\).

For general matrices, the geometry and coalescence of pseudospectral
components have been studied by variational and distance-to-defectivity
methods \cite{AlamEtAl2011,BurkeLewisOverton2003,LewisPang2008}.  The
present work complements that theory by determining, for the clean
scalar tridiagonal Toeplitz path, an exact finite-dimensional
connectedness criterion, contractibility of every component, and a
strict comparison of the merger barrier in consecutive dimensions.

Our contributions are as follows.  First, we prove vertical monotonicity
of the least singular value, which reduces planar topology to a real-axis
barrier problem and shows, more strongly, that every pseudospectral
component is contractible.  Second, a signed persymmetric pencil and a
folded Chebyshev parametrization give the strict comparison of the
largest barrier in consecutive dimensions.  Third, rectangular
factorizations give sharp-order two-sided bounds, explicit Lambert-$W$
bounds for the comparison-envelope threshold, and the resulting asymptotic
finite-size law.  Connectedness concerns the union over all
\emph{unstructured} perturbations: it does not imply that any single
perturbation supplies eigenvalues throughout a connecting region, nor
does it describe the attainable set for a prescribed structured
perturbation family.  Its physical content is instead the loss of a
component-wise, perturbation-independent spectral label at the prescribed
resolution.  The transition is also not the onset of the skin effect, which
is already present in the dimensionless gauge similarity above.

The paper is organized as follows.  \Cref{sec:main-results} states the
main result.  Vertical monotonicity is proved in
\cref{sec:vertical}; the strict barrier comparison occupies
\cref{sec:mesh}; and the topological criterion and quantitative bounds
are established in \cref{sec:topology-bounds}.  Numerical illustrations
and the physical interpretation are given in
\cref{sec:numerics,sec:physical-conclusion}.

\section{Model and main results}\label{sec:main-results}

For $n\geq2$ and $0<a<1$, let
\[
 A_n(a)=\tridiag(a,0,1).
\]
When \(a\) is fixed, we write \(A_n\) for \(A_n(a)\).
For $\varepsilon>0$, write
\[
 \sigma_\varepsilon(A_n)=
 \{z\in\C:s_{\min}(zI_n-A_n)<\varepsilon\}.
\]
Define
\[
 N_e(a,\varepsilon)=\min\{N\geq2:
 \sigma_\varepsilon(A_n)\text{ is connected for every }n\geq N\},
\]
\[
 N_f(a,\varepsilon)=\min\{N\geq2:
 \sigma_\varepsilon(A_N)\text{ is connected}\}.
\]
The following theorem proves that both defining sets are nonempty.

For the remainder of the article, put
\[
 r=\sqrt a,\qquad
 J_n=[\boldsymbol 1_{\{i+j=n+1\}}]_{i,j=1}^n.
\]
Thus \(A_n=V_n+aV_n^T\).  Here, \(\mathbf e_j\) denotes a Euclidean
coordinate vector of the ambient dimension, and \(U_j(t)\) denotes the
Chebyshev polynomial of the second kind, with \(U_{-1}=0\).

When the dimension is \(n\), write \(N=n+1\) and define
\begin{equation}\label{eq:main-definitions}
 I_n^{\mathrm{sp}}=
 \left[-2r\cos\frac{\pi}{N},2r\cos\frac{\pi}{N}\right],
 \quad
 g_n(x)=s_{\min}(xI_n-A_n),
 \quad
 \gamma_n=\max_{x\in I_n^{\mathrm{sp}}}g_n(x).
\end{equation}
The explicit lower and upper barriers used below are
\begin{align}
 \mathcal U_n(a)&=r^n\csc\frac{\pi}{N},\label{eq:main-U-def}\\
 \mathcal L_n(a)&=
 \left(1+r^2-2r\cos\frac{\pi}{N}\right)
 \frac{r^n(1-r^2)}
 {(1-r^{2N})\sin\frac{3\pi}{2N}}.\label{eq:main-L-def}
\end{align}

\begin{theorem}[Main theorem]\label{thm:main}
For every $0<a<1$ and $\varepsilon>0$, the following assertions hold.
\begin{enumerate}
\item For every $n\geq2$, every connected component of
$\sigma_\varepsilon(A_n(a))$ is
contractible.
\item For $n\geq2$, the barrier sequence is strictly decreasing,
$\gamma_{n+1}<\gamma_n$, tends to zero, and
\[
 \sigma_\varepsilon(A_n(a))\text{ is connected}
 \quad\Longleftrightarrow\quad \varepsilon>\gamma_n.
\]
\item The first-connectedness and eventual thresholds coincide:
\[
 N_c(a,\varepsilon):=N_f(a,\varepsilon)=N_e(a,\varepsilon).
\]
Moreover,
\begin{equation}\label{eq:N-two-sided-main}
 1+\max\{n\geq2:\mathcal L_n(a)\geq\varepsilon\}
 \leq N_c(a,\varepsilon)\leq
 \min\{n\geq2:\mathcal U_n(a)<\varepsilon\},
\end{equation}
where the maximum of the empty set is understood to be $1$.
\item For fixed $a$, as $\varepsilon\downarrow0$,
\begin{equation}\label{eq:N-asymptotic-main}
 N_c(a,\varepsilon)=\frac{2}{\log(1/a)}
 \left(\log(1/\varepsilon)+\log\!\log(1/\varepsilon)\right)+O_a(1),
\end{equation}
and $N_c(a,\varepsilon)=2$ if and only if $\varepsilon>a$.
\end{enumerate}
\end{theorem}

\section{Vertical monotonicity}\label{sec:vertical}

We first prove the vertical monotonicity that will be used in the
topological part of the argument.

\begin{lemma}[Vertical monotonicity]\label{lem:vertical}
For every real \(x\),
\[
 y\longmapsto s_{\min}((x+\ii y)I_n-A_n),\qquad y\geq0,
\]
is nondecreasing.
\end{lemma}

\begin{proof}
Put \(Y=y^2\), and, for \(t\geq0\), define
\[
 \mathcal P_n(Y,t)=((x+\ii\sqrt Y)I_n-A_n)^*
 ((x+\ii\sqrt Y)I_n-A_n)-t^2I_n,
 \qquad F_n(Y,t)=\det \mathcal P_n(Y,t).
\]
The same definition is used in every dimension: for \(k\geq1\),
\(\mathcal P_k(Y,t)\) is obtained by replacing \(n\) by \(k\),
\(F_k(Y,t)=\det \mathcal P_k(Y,t)\), and \(F_0=1\).
We shall prove the finite-dimensional implication
\begin{equation}\label{eq:det-derivative-positive}
 \mathcal P_n(Y,t)>0\quad\Longrightarrow\quad
 \partial_YF_n(Y,t)>0.
\end{equation}

At a fixed value of \(Y\), write
\[
 d=x^2+Y+1+a^2-t^2,
 \qquad b=-(1+a)x+\ii(1-a)\sqrt Y.
\]
For \(n\geq2\), the matrix \(\mathcal P_n(Y,t)\) has diagonal
\[
 d-1,d,\ldots,d,d-a^2,
\]
whereas for \(n=1\) its sole entry is \(d-1-a^2\).  In all dimensions
its first upper diagonal is \(b\), its first lower diagonal is \(\bar b\), and its second off-diagonals are \(a\).  Let \(T_k\) be the determinant of the pure
pentadiagonal Toeplitz matrix with diagonal \(d\) and these two
off-diagonals.  Set \(T_{-1}=0,T_0=1\), and put
\begin{equation}\label{eq:D-R-def}
 D_k=T_k-T_{k-1},\qquad R_k=T_k-a^2T_{k-1}.
\end{equation}
Indeed, a determinant is affine in each diagonal entry.  Replacing the
first diagonal entry \(d\) by \(d-1\) subtracts the cofactor
\(T_{k-1}\), whereas replacing the last one by \(d-a^2\) subtracts
\(a^2T_{k-1}\).  Thus, for \(k<n\), \(D_k\) and \(R_k\) are,
respectively, the leading and trailing principal minors of
\(\mathcal P_n(Y,t)\).

We record the continuant calculation that contains all the dependence on
\(Y\).  With \(T_k=0\) for \(k<0\), a Laplace expansion over the last
two rows, grouped according to the largest displaced index in the
contributing permutation, gives for every \(k\geq1\)
\begin{equation}\label{eq:T-recurrence}
\begin{aligned}
T_k={}&dT_{k-1}-(|b|^2-a^2)T_{k-2}
-\bigl(2a^2d-a(b^2+\bar b^2)\bigr)T_{k-3}\\
&-a^2(|b|^2-a^2)T_{k-4}+a^4dT_{k-5}-a^6T_{k-6}
-a^2\boldsymbol 1_{\{k=2\}}.
\end{aligned}
\end{equation}
For example, \(T_1=d\) and \(T_2=d^2-|b|^2\); these also verify the
only inhomogeneous coefficient in \eqref{eq:T-recurrence}.  Multiplying
the recurrence by \(z^k\) and summing gives
\begin{equation}\label{eq:T-generating}
 \sum_{k\geq0}T_kz^k=\frac{1-a^2z^2}{Q(z,Y)},
\end{equation}
where
\[
\begin{aligned}
 Q(z,Y)={}&1-dz+(|b|^2-a^2)z^2
 +\bigl(2a^2d-a(b^2+\bar b^2)\bigr)z^3\\
 &+a^2(|b|^2-a^2)z^4-a^4dz^5+a^6z^6.
\end{aligned}
\]
Consequently, with \(\mathcal D(z,Y)=\sum_{k\geq0}D_kz^k\),
\begin{equation}\label{eq:D-F-generating}
 \mathcal D(z,Y)=\frac{(1-z)(1-a^2z^2)}{Q(z,Y)},\qquad
 \sum_{k\geq0}F_k(Y,t)z^k=(1-a^2z)\mathcal D(z,Y).
\end{equation}
For the second identity, first change the initial diagonal, producing
\(D_k\), and then change the final diagonal; its cofactor is
\(D_{k-1}\).  Thus \(F_k=D_k-a^2D_{k-1}\), which proves the asserted
generating function.  Direct substitution of \(d\) and \(b\) gives
\[
 \partial_YQ=-z(1-z)(1+az)^2(1-a^2z).
\]
Differentiating \eqref{eq:D-F-generating} therefore yields
\begin{equation}\label{eq:F-prime-generating}
 \sum_{k\geq0}\partial_YF_k(Y,t)z^k
 =\varphi_a(z)\mathcal D(z,Y)^2,
 \qquad
 \varphi_a(z)=\frac{z(1+az)(1-a^2z)^2}{1-az}.
\end{equation}

Suppose that \(\mathcal P_n(Y,t)>0\).  Then
\begin{equation}\label{eq:D-R-positive}
 D_k>0,\qquad R_k>0\qquad(0\leq k<n).
\end{equation}
Put \(u_k=\sum_{j=0}^kD_jD_{k-j}\), with \(u_k=0\) for \(k<0\).
The coefficients of \(\varphi_a\), beginning with the coefficient of
\(z\),
are
\[
 1,\quad 2a(1-a),\quad a^2(a^2-4a+2),\quad
 2a^{k-1}(1-a)^2\quad(k\geq4).
\]
Thus
\begin{equation}\label{eq:F-prime-convolution}
\begin{aligned}
 \partial_YF_n={}&u_{n-1}+2a(1-a)u_{n-2}
 +a^2(a^2-4a+2)u_{n-3}\\
 &+2(1-a)^2\sum_{k=4}^na^{k-1}u_{n-k}.
\end{aligned}
\end{equation}

Only the third coefficient can be negative.  We now give the cofactor
inequality that absorbs it.  Let \(j,k\geq1\), \(j+k=n-1\).  Deleting
index \(k+1\) from \(\mathcal P_n\) leaves blocks of sizes \(k\) and \(j\),
coupled only by the second off-diagonal \(a\).  In block form, the
deleted principal submatrix is
\[
 \begin{bmatrix}
  L_k&a \mathbf e_k\mathbf e_1^T\\
  a \mathbf e_1\mathbf e_k^T&R_j^{\rm mat}
 \end{bmatrix},
 \qquad \det L_k=D_k,\quad \det R_j^{\rm mat}=R_j.
\]
Expanding once in each of the two coupling entries shows that its
determinant is
\begin{equation}\label{eq:deleted-cofactor}
 D_kR_j-a^2D_{k-1}R_{j-1}>0.
\end{equation}
Deleting index \(j+1\) instead gives the same inequality with \(j,k\)
interchanged.  Define
\[
 \vartheta_k=\frac{D_k}{D_{k-1}},\qquad
 \varpi_k=\frac{R_k}{R_{k-1}}.
\]
Since \(T_k=\sum_{\ell=0}^kD_\ell\), \eqref{eq:D-R-def} gives
\begin{equation}\label{eq:R-D-relation}
 R_k=D_k+(1-a^2)\sum_{\ell=0}^{k-1}D_\ell,\qquad
 R_k-R_{k-1}=D_k-a^2D_{k-1}.
\end{equation}
Hence \(\vartheta_k\leq a^2\) implies \(\varpi_k\leq1\).  On dividing
\eqref{eq:deleted-cofactor} and its companion by positive factors, we
obtain
\begin{equation}\label{eq:ratio-cross}
 \vartheta_k\varpi_j>a^2,\qquad \vartheta_j\varpi_k>a^2.
\end{equation}
If both \(\vartheta_j\) and \(\vartheta_k\) were at most \(a^2\), then
both \(\varpi_j\) and \(\varpi_k\) would be at most \(1\),
contradicting \eqref{eq:ratio-cross}.  Hence, after interchanging
\(j,k\) if necessary, \(\vartheta_j>a^2\).  Moreover,
\[
 \vartheta_jR_{j-1}-a^2R_j
 =(\vartheta_j-a^2)(R_{j-1}-a^2D_{j-1}),
\]
and the second factor equals
\((1-a^2)\sum_{\ell=0}^{j-1}D_\ell>0\).  Therefore
\[
 \vartheta_j>a^2\quad\Longrightarrow\quad
 \varpi_j<\vartheta_j/a^2.
\]
The first inequality in \eqref{eq:ratio-cross} now gives
\(a^2<\vartheta_k\varpi_j<\vartheta_j\vartheta_k/a^2\), and therefore
\begin{equation}\label{eq:D-ratio-product}
 \vartheta_j\vartheta_k>a^4.
\end{equation}
Consequently, for \(n\geq3\),
applying \eqref{eq:D-ratio-product} to \(j=p+1\), \(k=n-2-p\) gives
\[
 D_{p+1}D_{n-2-p}>a^4D_pD_{n-3-p},
 \qquad0\leq p\leq n-3.
\]
Summing and then adding the two positive endpoint terms gives
\begin{equation}\label{eq:u-absorption}
 u_{n-1}=2D_{n-1}+\sum_{j=1}^{n-2}D_jD_{n-1-j}>a^4u_{n-3}.
\end{equation}

If the third coefficient in \eqref{eq:F-prime-convolution} is negative,
write it as \(-\mathcal B_{\rm neg}\), where
\(\mathcal B_{\rm neg}=a^2(4a-2-a^2)\).  Since
\[
 a^4-\mathcal B_{\rm neg}=2a^2(1-a)^2>0,
\]
equations \eqref{eq:F-prime-convolution} and \eqref{eq:u-absorption}
prove \eqref{eq:det-derivative-positive}.  If that coefficient is
nonnegative, the implication is immediate.  For \(n=1\), one has
\(\partial_YF_1=1\); for \(n=2\), one has
\(\partial_YF_2=2D_1+2a(1-a)>0\).  Thus
\eqref{eq:det-derivative-positive} holds in every dimension.

Finally fix \(Y_0\geq0\).  If the least singular value at
\(Y_0\) is zero, the desired inequality is automatic.  Otherwise take
\(0\leq t<s_{\min}((x+\ii\sqrt{Y_0})I_n-A_n)\).  On the maximal interval to the
right of \(Y_0\) on which \(\mathcal P_n(Y,t)>0\), implication
\eqref{eq:det-derivative-positive} says that its determinant is strictly
increasing.  By continuity of Hermitian eigenvalues, a finite right
endpoint would be a first loss of positive definiteness at a
positive-semidefinite singular matrix, and hence would have determinant
zero.  This is impossible because the determinant has increased from
its positive value at \(Y_0\).
Thus \(\mathcal P_n(Y,t)>0\) for every \(Y\geq Y_0\).  Letting \(t\) increase to
the least singular value at \(Y_0\) proves the lemma.
\end{proof}

\section{Strict finite-size comparison of real-gap barriers}
\label{sec:mesh}

We next establish the strict comparison between consecutive dimensions.
Write the eigenvalues in decreasing order as
\[
 \lambda_{n,k}=2r\cos\frac{k\pi}{n+1},\qquad 1\leq k\leq n,
\]
and call
\[
 G_{n,k}=[\lambda_{n,k+1},\lambda_{n,k}],\qquad
 \gamma_{n,k}=\max_{x\in G_{n,k}}g_n(x),\qquad 1\leq k<n,
\]
the closed spectral gaps and their heights.  Continuity gives existence of
every maximum, and the spectral endpoints have height zero.  Every interior
gap point lies outside the spectrum, so \(xI_n-A_n\) is nonsingular there and
\(g_n(x)>0\).  Consequently every gap height is positive and every point at
which it is attained lies in the open gap.  Hence
\begin{equation}\label{eq:gamma-def}
 \gamma_n=\max_{1\leq k<n}\gamma_{n,k}.
\end{equation}

\begin{lemma}[Strict barrier comparison]\label{lem:mesh}
For every \(n\geq2\),
\[
 \gamma_{n+1}<\gamma_n.
\]
\end{lemma}

We give the complete finite-path calculation.  Persymmetry gives
\(A_nJ_n=J_nA_n^T\), so
\(B_n(x)=(xI_n-A_n)J_n\) is real symmetric.  Moreover,
\[
 (xI_n-A_n-sJ_n)J_n=B_n(x)-sI_n.
\]
Since \(J_n\) is orthogonal, the absolute eigenvalues of \(B_n(x)\)
are exactly the singular values of \(xI_n-A_n\), including multiplicity.
Consequently,
\begin{equation}\label{eq:signed-pencil}
 s_{\min}(xI_n-A_n)
 =\min\{|s|:\det(xI_n-A_n-sJ_n)=0\}.
\end{equation}
Let \(q_n(x,s)=\det(xI_n-A_n-sJ_n)\).
\begin{lemma}[The middle signed branch]\label{lem:middle-branch}
Write the eigenvalues of the real symmetric matrix
\[
 B_n(x)=(xI_n-A_n)J_n
\]
in decreasing order as
\(\alpha_1(x)\geq\cdots\geq\alpha_n(x)\), and abbreviate
\(\lambda_k=\lambda_{n,k}\).
If \(\lambda_{j+1}<x<\lambda_j\), then
\[
 s_{\min}(xI_n-A_n)=|s_*(x)|,
 \qquad \operatorname{sgn}(s_*(x))=(-1)^j,
\]
where
\[
 s_*(x)=
 \begin{cases}
  \alpha_m(x),&n=2m,\\
  \alpha_{m+1}(x),&n=2m+1.
 \end{cases}
\]
\end{lemma}

\begin{proof}
Let \(A_n^{\rm sym}=r(V_n+V_n^T)\).  Since
\(A_n=\mathcal G_n(r)A_n^{\rm sym}\mathcal G_n(r)^{-1}\) and
\(\mathcal G_n(r)^{-1}J_n
=r^{-(n-1)}J_n\mathcal G_n(r)\),
\begin{equation}\label{eq:B-congruence}
 B_n(x)=r^{-(n-1)}\mathcal G_n(r)(xI_n-A_n^{\rm sym})
 J_n\mathcal G_n(r).
\end{equation}
The sine eigenvector of \(A_n^{\rm sym}\) belonging to \(\lambda_k\) is an
eigenvector of \(J_n\) with eigenvalue \((-1)^{k+1}\).  Hence
\eqref{eq:B-congruence} and Sylvester's law of inertia show that
\(B_n(x)\) has the same inertia as the diagonal matrix with entries
\[
 (-1)^{k+1}(x-\lambda_k),\qquad 1\leq k\leq n.
\]
Counting signs gives
\begin{equation}\label{eq:middle-inertia}
 \nu_+(B_{2m}(x))=
 \begin{cases}m,&j\text{ even},\\m-1,&j\text{ odd},\end{cases}
 \qquad
 \nu_+(B_{2m+1}(x))=
 \begin{cases}m+1,&j\text{ even},\\m,&j\text{ odd}.
 \end{cases}
\end{equation}

First let \(n=2m\).  Permute the coordinates into the two groups
\[
 (1,\ldots,m),\qquad(2m,2m-1,\ldots,m+1).
\]
In these coordinates
\begin{equation}\label{eq:even-dilation}
\begin{aligned}
 B_{2m}(x)&=\widehat P_m(x)-\mathcal R_{\rm bd},\\
 \widehat P_m(x)&=\begin{bmatrix}0&P_m(x)\\P_m(x)^T&0\end{bmatrix},\\
 \mathcal R_{\rm bd}
 &=\begin{bmatrix}\mathbf e_m\mathbf e_m^T&0\\
 0&a \mathbf e_m\mathbf e_m^T\end{bmatrix},
\end{aligned}
\end{equation}
where \(P_m(x)=xI_m-A_m\).  Thus \(\mathcal R_{\rm bd}\geq0\) and
\(\operatorname{rank}\mathcal R_{\rm bd}=2\).  For
\(\Gamma=\diag(I_m,-I_m)\),
\[
 \widetilde B=-\Gamma B_{2m}(x)\Gamma
 =\widehat P_m(x)+\mathcal R_{\rm bd}
 =B_{2m}(x)+2\mathcal R_{\rm bd}.
\]
If \(\widetilde\alpha_1\geq\cdots\geq\widetilde\alpha_{2m}\) are the
eigenvalues of
\(\widetilde B\), then
\begin{equation}\label{eq:alpha-reflection}
 \widetilde\alpha_k=-\alpha_{2m+1-k}.
\end{equation}
Positive-semidefinite rank-two interlacing gives
\begin{equation}\label{eq:rank-two-interlace}
 \widetilde\alpha_k\geq\alpha_k\geq\widetilde\alpha_{k+2},
 \qquad1\leq k\leq2m-2.
\end{equation}
If \(j\) is even, the only candidates for minimum modulus
have moduli \(\alpha_m\) and
\(-\alpha_{m+1}=\widetilde\alpha_m\), and
\(\alpha_m\leq\widetilde\alpha_m\).  If \(j\) is odd and \(m\geq2\), the two
candidates have moduli \(\alpha_{m-1}\) and
\(-\alpha_m=\widetilde\alpha_{m+1}\); taking \(k=m-1\) in
\eqref{eq:rank-two-interlace} gives
\(\alpha_{m-1}\geq\widetilde\alpha_{m+1}\).  For \(m=1\), both eigenvalues in
the sole gap are negative, so the larger one \(\alpha_1\) has minimum
modulus.  Thus \(\alpha_m\) always attains the minimum modulus, and its
sign follows from \eqref{eq:middle-inertia}.  Equality in an interlacing
comparison is harmless: it only gives a multiple least singular value.

Now let \(n=2m+1\).  Group the first and last \(m\) coordinates as
above and leave the middle coordinate last.  Then
\begin{equation}\label{eq:odd-dilation}
 B_{2m+1}(x)=
 \begin{bmatrix}\widehat P_m(x)&\boldsymbol w\\\boldsymbol w^T&x\end{bmatrix},
 \qquad
 \boldsymbol w=\begin{bmatrix}-\mathbf e_m\\-a\mathbf e_m\end{bmatrix}.
\end{equation}
In an open spectral gap of \(A_{2m+1}\), \(P_m(x)\) is nonsingular.  Indeed,
the explicit path spectrum gives
\[
 \operatorname{spec}(A_m)
 =\left\{2r\cos\frac{\ell\pi}{m+1}:1\leq\ell\leq m\right\}
 =\left\{\lambda_{2m+1,2\ell}:1\leq\ell\leq m\right\},
\]
so its zeros are exactly the even-indexed eigenvalues of \(A_{2m+1}\), none
of which lies in an open spectral gap.
The eigenvalues of the principal submatrix \(\widehat P_m(x)\) are
consequently
\[
 s_1(P_m)\geq\cdots\geq s_m(P_m)>0>
 -s_m(P_m)\geq\cdots\geq-s_1(P_m).
\]
Cauchy interlacing gives
\[
 \alpha_m\geq s_m(P_m)\geq\alpha_{m+1}
 \geq-s_m(P_m)\geq\alpha_{m+2}.
\]
If \(j\) is even, \eqref{eq:middle-inertia} and interlacing give
\[
 0<\alpha_{m+1}\leq s_m(P_m)\leq-\alpha_{m+2};
\]
if \(j\) is odd, they give
\[
 0<-\alpha_{m+1}\leq s_m(P_m)\leq\alpha_m.
\]
Every other eigenvalue is farther from zero.  Hence
\(\alpha_{m+1}\) attains the minimum modulus, with sign prescribed by
\eqref{eq:middle-inertia}; equality again only means multiplicity.  At a
spectral endpoint the zero is simple because \(A_n\) is similar to the
Jacobi matrix \(A_n^{\rm sym}\), whose eigenvalues are simple.  Thus the selected
ordered branch extends continuously to zero at both ends of the closed
gap.  Finally, multiplication by \(J_n\) does not change singular
values, which proves the assertion.
\end{proof}

\begin{proof}[Proof of \cref{lem:mesh}]
The proof proceeds in five stages: (i) derivation of parity-unified
folded-continuant formulas; (ii) conversion to the half-angle chord
parametrization; (iii) comparison of all noncentral gaps in consecutive
dimensions; (iv) separate treatment of the central gaps by interlacing; and
(v) assembly of the global strict barrier inequality.  For orientation,
\(F_n\) denotes the determinant from the vertical argument, whereas
\(\mathcal F_h\) below is an angular factor; the decorated
\(\mathcal C,\mathcal D,\mathcal R_{\rm bd}\) are local objects; and
\(\mathfrak c_h,\mathfrak d_h\) are hyperbolic or angular quantities, unlike
the central gap heights \(c_m,d_m\).

\medskip\noindent\emph{Folded continuants.}\quad
Put
\[
 C_0=x^2+(1-a)^2-s^2,\qquad
 \zeta_s=\frac{1+a^2+(1+a)s}{2a},
\]
and define the folded discriminant and recurrence coefficient by
\[
 \Delta_{\rm fold}=((1+a-x)^2-s^2)((1+a+x)^2-s^2),
\]
\[
 \mathfrak b=(1+a^2)x^2+2as^2+2a^2-2a-2a^3.
\]
The global recurrence polynomial is
\begin{equation}\label{eq:fold-characteristic}
\begin{gathered}
 \Pi(z):=1-C_0z+\mathfrak b z^2-a^2C_0z^3+a^4z^4,\\
 \left.
 \begin{aligned}
  2a(u+v)&=C_0,\\
  a^2(4uv+2)&=\mathfrak b
 \end{aligned}
 \right\}
 \Longrightarrow
 \Pi(z)=(1-2auz+a^2z^2)(1-2avz+a^2z^2),
 \qquad u,v\in\mathbb R.
\end{gathered}
\end{equation}
On a real sheet with \(\Delta_{\rm fold}\geq0\), we additionally put
\[
 \xi_\pm=\frac{C_0\pm\sqrt{\Delta_{\rm fold}}}{4a}.
\]
Then the pair \((u,v)=(\xi_+,\xi_-)\) satisfies the two symmetric
relations in \eqref{eq:fold-characteristic}:
\[
 \xi_++\xi_-=\frac{C_0}{2a},\qquad
 \mathfrak b=a^2(4\xi_+\xi_-+2).
\]
The proof uses the root presentation only on these real sheets; the global
transfer and generating identities use the polynomial \(\Pi\) itself.
Here is a direct continuant derivation of the required folded formulas.
Pair the coordinates in the order
\((1,n),(2,n-1),\ldots\).  In this order the symmetric matrix
\((xI_n-A_n)J_n-sI_n\) is block tridiagonal.  Its interior diagonal and
off-diagonal blocks are
\[
 D_{\rm int}=\begin{bmatrix}-s&x\\x&-s\end{bmatrix},
 \qquad
 E_{\rm off}=\begin{bmatrix}0&-1\\-a&0\end{bmatrix},
\]
with lower block \(E_{\rm off}^T\).  For even order the terminal block is
\[
 D_{\rm term}^{\rm e}=\begin{bmatrix}-1-s&x\\x&-a-s\end{bmatrix};
\]
for odd order the terminal block is the scalar \(x-s\), coupled to the
preceding block by the column \((-1,-a)^T\).  Since
\[
 (xI_n-A_n-sJ_n)=\bigl((xI_n-A_n)J_n-sI_n\bigr)J_n,
\]
the determinant of this block matrix, multiplied by \(\det J_n\), is
\(q_n(x,s)\).  Set \(\mathcal Q_m^{\rm e}=q_{2m}(x,s)\),
\(\mathcal Q_m^{\rm o}=q_{2m+1}(x,s)\).  We include the continuant
verification to fix all boundary signs.
We now give a complete minor transfer for the folded continuant.
For \(\tau\in\{\mathrm e,\mathrm o\}\), let \(K_m^\tau\) denote the
corresponding matrix
\[
 (xI_n-A_n-sJ_n)J_n=(xI_n-A_n)J_n-sI_n,
 \qquad n=2m\ \hbox{or}\ 2m+1,
\]
after simultaneously permuting its rows and columns into the paired
order.  Thus \(K_m^{\mathrm e}\) has order \(2m\),
\(K_m^{\mathrm o}\) has order \(2m+1\), and
\[
 K_1^{\mathrm e}=
 \begin{bmatrix}-1-s&x\\x&-a-s\end{bmatrix},
 \qquad
 K_1^{\mathrm o}=
 \begin{bmatrix}
 -s&x&-1\\
 x&-s&-a\\
 -1&-a&x-s
 \end{bmatrix}.
\]
For \(m\geq1\), increasing \(m\) by one adds one interior block at the
exposed end:
\begin{equation}\label{eq:fold-block-extension}
 K_{m+1}^\tau=
 \begin{bmatrix}
  D_{\rm int}&\mathcal E_m^\tau\\
  (\mathcal E_m^\tau)^T&K_m^\tau
 \end{bmatrix},
 \qquad
 \mathcal E_m^\tau=
 \begin{bmatrix}E_{\rm off}&0\end{bmatrix},
\end{equation}
where the zero block has the width required by \(K_m^\tau\).

For a matrix whose first two coordinates are exposed, write
\(K(\widehat I\mid\widehat J)\) for the matrix obtained by deleting
rows \(I\) and columns \(J\).  Suppressing the parity superscript, set
\[
\begin{aligned}
 \mathfrak d_m&=\det K_m,&
 \mathfrak p_m&=\det K_m(\widehat1\mid\widehat1),&
 \mathfrak q_m&=\det K_m(\widehat2\mid\widehat2),\\
 \mathfrak c_m&=\det K_m(\widehat1\mid\widehat2),&
 \mathfrak r_m&=\det K_m(\widehat{1,2}\mid\widehat{1,2}).
\end{aligned}
\]
The off-principal minor obtained by deleting row \(2\) and column \(1\)
has the same value \(\mathfrak c_m\), because \(K_m\) is symmetric.

Here is the one-step calculation.  On the dense set on which both
\(K_m\) and the tail obtained by deleting its first two coordinates
are invertible, let
\[
 F=\begin{bmatrix}u&w\\w&v\end{bmatrix}
\]
be its Schur complement in \(K_m\), and let \(\delta\) be the
determinant of that tail.  Then
\[
 (\mathfrak d_m,\mathfrak p_m,\mathfrak q_m,\mathfrak c_m,\mathfrak r_m)
 =\delta(uv-w^2,v,u,w,1).
\]
Since
\[
 E_{\rm off}\operatorname{adj}(F)E_{\rm off}^T
 =
 \begin{bmatrix}u&-aw\\-aw&a^2v\end{bmatrix},
\]
taking the Schur complement of \(K_m\) in
\eqref{eq:fold-block-extension} gives
\begin{equation}\label{eq:five-minor-update}
\begin{aligned}
 \mathfrak d_{m+1}
 &= (s^2-x^2)\mathfrak d_m
    +s(a^2\mathfrak p_m+\mathfrak q_m)
    -2ax\mathfrak c_m+a^2\mathfrak r_m,\\
 \mathfrak p_{m+1}&=-s\mathfrak d_m-a^2\mathfrak p_m,\\
 \mathfrak q_{m+1}&=-s\mathfrak d_m-\mathfrak q_m,\\
 \mathfrak c_{m+1}&=x\mathfrak d_m+a\mathfrak c_m,\\
 \mathfrak r_{m+1}&=\mathfrak d_m.
\end{aligned}
\end{equation}
Both sides are polynomial in \(x\) and \(s\), and the nonsingular set used
above is dense.  The polynomial identity theorem therefore extends all five
identities to the singular cases as well; no inverse is assumed in
\eqref{eq:five-minor-update}.

The five quantities in \eqref{eq:five-minor-update} satisfy one linear
invariant.  Define
\[
 \ell_m=(1-a)\mathfrak d_m-as\,\mathfrak p_m+s\,\mathfrak q_m
       +(1-a)x\mathfrak c_m-a(1-a)\mathfrak r_m.
\]
Direct substitution into \eqref{eq:five-minor-update} gives
\[
 \ell_{m+1}=-a\ell_m.
\]
For the two terminal blocks, direct evaluation gives
\[
\begin{aligned}
&(\mathfrak d_1,\mathfrak p_1,\mathfrak q_1,\mathfrak c_1,\mathfrak r_1)_{\mathrm e}\\
&\qquad=\bigl((1+s)(a+s)-x^2,-a-s,-1-s,x,1\bigr),\\
&(\mathfrak d_1,\mathfrak p_1,\mathfrak q_1,\mathfrak c_1,\mathfrak r_1)_{\mathrm o}\\
&\qquad=\bigl(x(1+a^2)+2as-C_0(x-s),\\
&\hspace{26mm}s^2-sx-a^2,\ s^2-sx-1,\ x^2-sx-a,\ x-s\bigr).
\end{aligned}
\]
Thus \(\ell_1=0\) in both parity classes, and consequently
\begin{equation}\label{eq:fold-minor-invariant}
 (1-a)\mathfrak d_m-as\,\mathfrak p_m+s\,\mathfrak q_m
 +(1-a)x\mathfrak c_m-a(1-a)\mathfrak r_m=0
 \qquad(m\geq1).
\end{equation}

Because
\[
 \det J_{2m}=\det J_{2m+1}=(-1)^m,
\]
the first signed minor below is precisely the required determinant:
\[
 \boldsymbol\eta_m^\tau
 =
 (-1)^m
 \begin{bmatrix}
 \mathfrak d_m^\tau&\mathfrak p_m^\tau&
 \mathfrak q_m^\tau&\mathfrak c_m^\tau
 \end{bmatrix}^{T},
 \qquad
 (\boldsymbol\eta_m^{\mathrm e})_1=\mathcal Q_m^{\mathrm e},
 \quad
 (\boldsymbol\eta_m^{\mathrm o})_1=\mathcal Q_m^{\mathrm o}.
\]
Eliminating \(\mathfrak r_m\) with \eqref{eq:fold-minor-invariant} reduces
\eqref{eq:five-minor-update} to the four-dimensional transfer
\begin{equation}\label{eq:fold-four-transfer}
 \boldsymbol\eta_{m+1}^\tau
 =\mathsf T_{\rm fold}\boldsymbol\eta_m^\tau,
 \qquad
 \mathsf T_{\rm fold}=
 \begin{bmatrix}
 x^2-s^2-a&\dfrac{a^3s}{1-a}&-\dfrac{s}{1-a}&ax\\[2mm]
 s&a^2&0&0\\
 s&0&1&0\\
 -x&0&0&-a
 \end{bmatrix}.
\end{equation}
The terminal conditions can be incorporated at \(m=0\) by setting
\begin{equation}\label{eq:fold-transfer-initial}
 \boldsymbol\eta_0^{\mathrm e}
 =\begin{bmatrix}1&a^{-1}&1&0\end{bmatrix}^{T},
 \qquad
 \boldsymbol\eta_0^{\mathrm o}
 =\begin{bmatrix}x-s&1&1&-1\end{bmatrix}^{T}.
\end{equation}
Indeed, direct multiplication in \eqref{eq:fold-four-transfer} gives
\[
 \mathsf T_{\rm fold}\boldsymbol\eta_0^{\mathrm e}
 =
 \begin{bmatrix}
 x^2-(1+s)(a+s)\\ a+s\\1+s\\-x
 \end{bmatrix},
\]
and
\[
 \mathsf T_{\rm fold}\boldsymbol\eta_0^{\mathrm o}
 =
 \begin{bmatrix}
 C_0(x-s)-x(1+a^2)-2as\\
 a^2-s^2+sx\\
 1-s^2+sx\\
 a+sx-x^2
 \end{bmatrix},
\]
which are exactly the signed terminal minors above.  Hence
\(\boldsymbol\eta_m^\tau
=\mathsf T_{\rm fold}^{\,m}\boldsymbol\eta_0^\tau\) for every
\(m\geq0\).

A direct determinant calculation gives
\begin{equation}\label{eq:fold-transfer-characteristic}
\begin{aligned}
 \det(\lambda I_4-\mathsf T_{\rm fold})
 &=
 \lambda^4-C_0\lambda^3+\mathfrak b\lambda^2
 -a^2C_0\lambda+a^4,\\
 \mathfrak b
 &=(1+a^2)x^2+2as^2+2a^2-2a-2a^3.
\end{aligned}
\end{equation}
Cayley--Hamilton, followed by multiplication by
\(\mathsf T_{\rm fold}^{\,m-4}\boldsymbol\eta_0^\tau\) and extraction
of the first coordinate, therefore yields, for
\(Z_m=\mathcal Q_m^{\mathrm e}\) and for
\(Z_m=\mathcal Q_m^{\mathrm o}\),
\begin{equation}\label{eq:fold-recurrence}
 Z_m-C_0Z_{m-1}+\mathfrak bZ_{m-2}
 -a^2C_0Z_{m-3}+a^4Z_{m-4}=0,
 \qquad m\geq4.
\end{equation}

Finally, applying \(\mathsf T_{\rm fold}\) three times to
\eqref{eq:fold-transfer-initial} gives the boundary residuals
\[
 \mathsf{Res}_0=Z_0,\quad
 \mathsf{Res}_1=Z_1-C_0Z_0,\quad
 \mathsf{Res}_2=Z_2-C_0Z_1+\mathfrak bZ_0,
\]
\[
 \mathsf{Res}_3
 =Z_3-C_0Z_2+\mathfrak bZ_1-a^2C_0Z_0,
\]
namely
\[
\begin{array}{ccccc}
 \toprule
 &\mathsf{Res}_0&\mathsf{Res}_1&\mathsf{Res}_2&\mathsf{Res}_3\\ \midrule
 Z_m=\mathcal Q_m^{\rm e}
 &1&a(1-2\zeta_s)&a^2(1-2\zeta_s)&a^3\\
 Z_m=\mathcal Q_m^{\rm o}
 &x-s&-[x(1+a^2)+2as]&a^2(x-s)&0\\
 \bottomrule
\end{array}
\]
as asserted.
Thus, as a formal-power-series identity,
\[
 \Pi(z)\sum_{m\geq0}Z_mz^m
 =\mathsf{Res}_0+\mathsf{Res}_1z+\mathsf{Res}_2z^2+\mathsf{Res}_3z^3.
\]
For the even sequence the polynomial on the right factors as
\((1+az)(1-2a\zeta_sz+a^2z^2)\); for the odd sequence it is
\((x-s)-[x(1+a^2)+2as]z+a^2(x-s)z^2\).  This proves
\eqref{eq:fold-recurrence}.  Since \(\Pi(0)=1\), it is invertible in
\(\mathbb R[[z]]\), and the two generating functions are the literal
formal-power-series identities
\begin{equation}\label{eq:fold-generating-even}
 \sum_{m\geq0}\mathcal Q_m^{\rm e}z^m
 =(1+az)(1-2a\zeta_sz+a^2z^2)\Pi(z)^{-1},
\end{equation}
and
\begin{equation}\label{eq:fold-generating-odd}
 \sum_{m\geq0}\mathcal Q_m^{\rm o}z^m
 =\bigl((x-s)-[x(1+a^2)+2as]z+a^2(x-s)z^2\bigr)
  \Pi(z)^{-1}.
\end{equation}
Now let \(U_{-1}=0\).  The Chebyshev generating function gives
\[
 \sum_{m\geq0}(u-\zeta_s)(U_m(u)+U_{m-1}(u))\omega_{\rm gf}^m
 =\frac{(u-\zeta_s)(1+\omega_{\rm gf})}
 {1-2u\omega_{\rm gf}+\omega_{\rm gf}^2},
\]
and, with \(\mathsf B_s^{\rm fold}=[x(1+a^2)+2as]/(2a)\),
\[
 \sum_{m\geq0}((x-s)u-\mathsf B_s^{\rm fold})U_m(u)\omega_{\rm gf}^m
 =\frac{(x-s)u-\mathsf B_s^{\rm fold}}
 {1-2u\omega_{\rm gf}+\omega_{\rm gf}^2}.
\]
For a differentiable real function \(F\), write
\[
 [F]_{u,v}:=
 \begin{cases}
  \dfrac{F(u)-F(v)}{u-v},&u\neq v,\\[2mm]
  F'(u),&u=v.
 \end{cases}
\]
Let \(u,v\in\mathbb R\) satisfy the two symmetric relations in
\eqref{eq:fold-characteristic}.  Taking the corresponding divided
difference of the Chebyshev generating identities and putting
\(\omega_{\rm gf}=az\) recovers
\eqref{eq:fold-generating-even} and \eqref{eq:fold-generating-odd}.
Coefficient comparison gives
\begin{equation}\label{eq:fold-even}
 q_{2m}(x,s)=a^m[\mathfrak f_m]_{u,v},\qquad
 \mathfrak f_m(t)=(t-\zeta_s)(U_m(t)+U_{m-1}(t)),
\end{equation}
and
\begin{equation}\label{eq:fold-odd}
 q_{2m+1}(x,s)=a^m[\mathfrak g_m]_{u,v},\qquad
 \mathfrak g_m(t)=\left((x-s)t-
 \frac{x(1+a^2)+2as}{2a}\right)U_m(t).
\end{equation}
The determinant recurrence and generating identities are global.  The
real-pair formulas \eqref{eq:fold-even}--\eqref{eq:fold-odd} are used
only after the two symmetric relations have been verified; their
collision value is the displayed derivative.  No complex square-root
choice is made when \(\Delta_{\rm fold}<0\), and none is needed below.

\medskip\noindent\emph{Half-angle chord parametrization.}\quad
Recall that \(a=e^{-2h}\), and put
\[
 \mathfrak s_h=\sinh h,\qquad \mathfrak c_h=\cosh h.
\]
For \(K\geq2\) and \(0\leq\theta\leq\pi\), define
\begin{equation}\label{eq:Phi-Psi-def}
 \Phi_K(\theta)=\sqrt{\mathfrak s_h^2+\sin^2\theta}\,
 \frac{|\sin(K\theta)|}{\sin\theta},
\end{equation}
where the quotient is interpreted continuously at \(\theta=0,\pi\),
and
\[
 \Psi_K(\eta)=\sqrt{\mathfrak s_h^2-\sinh^2\eta}\,
 \frac{\sinh(K\eta)}{\sinh\eta},\qquad0\leq\eta\leq h,
\]
with the limiting value at \(\eta=0\).  We spell out the half-angle
reduction, including its sign.  Put
\[
 y=\cos\theta,\qquad z=\cosh\eta,\qquad
 \mathfrak p_\theta=\sqrt{\mathfrak c_h^2-y^2},\qquad
 \mathfrak p_\eta=\sqrt{\mathfrak c_h^2-z^2},
\]
and let \(L_0=2r/\mathfrak c_h\).  If
\(x=L_0yz\), \(s=L_0\mathfrak p_\theta\mathfrak p_\eta>0\), and
\(\tau=\varsigma s\) with \(\varsigma\in\{-1,1\}\), then direct
expansion gives
\begin{equation}\label{eq:half-variable-identities}
\begin{gathered}
 (1+a+x)^2-s^2=4a(z+y)^2,\qquad
 (1+a-x)^2-s^2=4a(z-y)^2,\\
 u=2z^2-1,\qquad v=2y^2-1,\qquad
 \zeta_\tau=2\mathfrak c_h^2-1
 +2\varsigma\mathfrak p_\theta\mathfrak p_\eta.
\end{gathered}
\end{equation}
The real pair \((u,v)\) satisfies the two symmetric relations in
\eqref{eq:fold-characteristic}.
For \(K=2m+1\), use
\[
 U_m(2t^2-1)+U_{m-1}(2t^2-1)=U_{K-1}(t);
\]
for \(K=2m+2\), use
\[
 U_m(2t^2-1)=\frac{U_{K-1}(t)}{2t}.
\]
Substitution of this pair in \eqref{eq:fold-even} and
\eqref{eq:fold-odd}, respectively, reduces either equation
\(q_{K-1}(x,\tau)=0\) to
\begin{equation}\label{eq:signed-half-chord}
 \mathfrak p_\theta U_{K-1}(y)
 =\varsigma\mathfrak p_\eta U_{K-1}(z).
\end{equation}
In the odd-fold case the reduction also contains the factor
\(\mathfrak p_\eta y+\varsigma\mathfrak p_\theta z\).  On a
distinct-variable sheet this factor cannot vanish: after squaring, its
vanishing would give
\[
 (\mathfrak c_h^2-z^2)y^2=(\mathfrak c_h^2-y^2)z^2,
\]
and hence \(y^2=z^2\).  Thus the factor is nonzero whenever the folded
variables are distinct.  On the outer lobe
\(U_{K-1}(z)>0\), whereas, for
\(j\pi/K<\theta<(j+1)\pi/K\),
\[
 \operatorname{sgn}\bigl(U_{K-1}(\cos\theta)\bigr)
 =\operatorname{sgn}\bigl(\sin(K\theta)\bigr)=(-1)^j.
\]
Moreover,
\(\mathfrak p_\theta|U_{K-1}(y)|=\Phi_K(\theta)\) and
\(\mathfrak p_\eta U_{K-1}(z)=\Psi_K(\eta)\).  Thus
\eqref{eq:signed-half-chord} is equivalent to
\(\varsigma=(-1)^j\) and
\begin{equation}\label{eq:chord-equation}
 \Phi_K(\theta)=\Psi_K(\eta),
\end{equation}
with
\begin{equation}\label{eq:channel-xs}
 s=\mathfrak s_{\rm hyp}(\theta,\eta)
 :=\frac{2r}{\mathfrak c_h}
 \sqrt{(\mathfrak s_h^2+\sin^2\theta)
       (\mathfrak s_h^2-\sinh^2\eta)},\qquad
 x=\frac{2r}{\mathfrak c_h}\cos\theta\cosh\eta.
\end{equation}
Consequently \(q_{K-1}(x,(-1)^js)=0\).  Conversely, for distinct real
folded variables, their sum and product determine
\[
 x^2=L_0^2y^2z^2,\qquad
 s^2=L_0^2(\mathfrak c_h^2-y^2)(\mathfrak c_h^2-z^2).
\]
Thus, for distinct real variables and a fixed sign of \(x\), the
preceding steps are reversible, and the two chord relations are
equivalent to the corresponding signed equation.  The determinant
identities themselves extend to a folded-variable collision by the
continuous divided-difference convention, but we invoke the chord
equivalence only for distinct variables; the selected middle branch is
proved below to be collision-free.
Here \(K-1\) is the matrix size.  If the exterior argument has crossed
through \(1\), write it as \(z=\cos\phi\),
\(0\leq\phi<\pi/K\).  Then the chord equation is
\(\Phi_K(\theta)=\Phi_K(\phi)\), and the second factor in the formula
for \(s\) is \(\mathfrak s_h^2+\sin^2\phi\).  We denote the resulting
value by
\[
 \mathfrak s_{\rm ell}(\theta,\phi)
 =\frac{2r}{\mathfrak c_h}
 \sqrt{(\mathfrak s_h^2+\sin^2\theta)
       (\mathfrak s_h^2+\sin^2\phi)}.
\]

\medskip\noindent\emph{Branch classification and exhaustion.}\quad
We next identify exactly which of these chords is the middle branch of
\Cref{lem:middle-branch}.  Set
\begin{equation}\label{eq:Wcal-def}
 \mathcal W_K(u)=\sqrt{\mathfrak c_h^2-u^2}\,|U_{K-1}(u)|,
 \qquad |u|\leq\mathfrak c_h.
\end{equation}
At \(u=\pm\mathfrak c_h\), this denotes the continuous endpoint value
\(0\).
Then
\[
 \Phi_K(\theta)=\mathcal W_K(\cos\theta),\qquad
 \Psi_K(\eta)=\mathcal W_K(\cosh\eta).
\]
If \(u_{\ell,K}^{(0)}=\cos\frac{\ell\pi}{K}\), then, away from the zeros,
\[
 \frac{d}{du}\log\mathcal W_K(u)
 =\sum_{\ell=1}^{K-1}\frac1{u-u_{\ell,K}^{(0)}}
 -\frac{u}{\mathfrak c_h^2-u^2},
\]
and its derivative is
\begin{equation}\label{eq:W-log-concave}
 -\sum_{\ell=1}^{K-1}\frac1{(u-u_{\ell,K}^{(0)})^2}
 -\frac{\mathfrak c_h^2+u^2}{(\mathfrak c_h^2-u^2)^2}<0.
\end{equation}
At each finite nodal endpoint the logarithmic derivative tends to the
appropriate signed infinity, and it tends to \(-\infty\) as
\(u\uparrow\mathfrak c_h\).  Thus there is exactly one lobe maximum in
every nodal
interval, including the outer interval
\(\mathcal J_K^{\rm out}=(\cos\frac{\pi}{K},\mathfrak c_h)\).  Let
\(u_K^{\max}\) be the unique maximizer on \(\mathcal J_K^{\rm out}\).
For every
\(0\leq\ell<\mathcal W_K(u_K^{\max})\), there is a unique solution
\[
 z_K(\ell)\in(u_K^{\max},\mathfrak c_h],\qquad
 \mathcal W_K(z_K(\ell))=\ell;
\]
we call it the \emph{endpoint-side outer solution}.  If
\(z_K(\ell)\geq1\), write it as \(\cosh\eta\),
\(0\leq\eta\leq h\); if it is at most \(1\), write it as
\(\cos\phi\), \(0\leq\phi<\pi/K\).  These are two coordinates on the
same outer lobe, joined at \(z=1\).

For a positive-half gap and
\(\theta\in(j\pi/K,(j+1)\pi/K)\), \(\theta\leq\pi/2\), a
\emph{$j$-chord} is the pair
\[
 y=\cos\theta,\qquad
 z=z_K(\mathcal W_K(y)),\qquad
 \mathcal W_K(y)=\mathcal W_K(z),
\]
together with the coordinate-free values
\[
 x=L_0yz,\qquad
 s=L_0\sqrt{(\mathfrak c_h^2-y^2)(\mathfrak c_h^2-z^2)};
\]
its signed root is
\((-1)^js\).  The term endpoint-side always refers to the part of the
outer lobe adjacent to \(\mathfrak c_h\), not to the hyperbolic or elliptic
coordinate used for it.

We now prove that these chords are exactly the middle signed branch.  By
reflection it suffices to consider a positive-half gap, so \(j\geq1\)
and \(\theta\leq\pi/2\).  Put
\(\phi_{\rm rem}=\theta-j\pi/K\in(0,\pi/K)\).  Since
\(|\sin(K\phi_{\rm rem})|=|\sin(K\theta)|\) and
\[
 \mathcal F_h(t)=
 \frac{\sqrt{\mathfrak s_h^2+\sin^2t}}{\sin t}
\]
is strictly decreasing on \((0,\pi/2]\),
\begin{equation}\label{eq:outer-level-dominates}
 \Phi_K(\theta)<\Phi_K(\phi_{\rm rem})
 \leq\max_{u\in\mathcal J_K^{\rm out}}\mathcal W_K(u).
\end{equation}
Thus the endpoint-side outer solution exists for every inner parameter.

At \(x=\lambda_{K-1,k}=2r\cos\frac{k\pi}{K}\), \(s=0\), the two
folded variables are \(\cos\frac{2k\pi}{K}\) and
\(\cosh(2h)\), or, in half variables, \(\cos\frac{k\pi}{K}\) and
\(\mathfrak c_h\).  Moreover, zero is a simple eigenvalue of
\(B_{K-1}(\lambda_{K-1,k})\): the matrix
\(\lambda_{K-1,k}I_{K-1}-A_{K-1}\) has rank \(K-2\), and \(J_{K-1}\) is
invertible.  The implicit-function theorem therefore gives, near either
gap endpoint, a real continuation of the middle signed root and real
half variables \(y=\cos\theta\) and \(z\).

Continue this ordered middle root \(s_*(x)\) through the open gap.  Write
\(\Delta_*(x)=\Delta_{\rm fold}(x,s_*(x))\) for its folded discriminant,
let \(y(x)\leq z(x)\) be the reconstructed nonnegative half variables on
the real sheet issuing from the endpoint germ, and set
\(\Theta(x)=\arccos y(x)\).  Let \(\mathcal C\) be the set of gap points for
which
\[
 \Delta_*(x)>0,\qquad
 \frac{j\pi}{K}<\Theta(x)<\frac{(j+1)\pi}{K},\qquad
 u_K^{\max}<z(x)<\mathfrak c_h.
\]
On \(\mathcal C\), the reversible folded determinant identities give
\(\mathcal W_K(y(x))=\mathcal W_K(z(x))\); uniqueness on the endpoint-side
outer lobe then gives
\[
 z(x)=z_K(\mathcal W_K(y(x))).
\]
Thus \(\mathcal C\) is precisely the selected endpoint-side chord sheet.
The endpoint germ just constructed makes \(\mathcal C\) nonempty.  The
ordered eigenvalues of the continuous real symmetric pencil \(B_{K-1}(x)\)
are continuous in \(x\), including at an eigenvalue crossing, so a multiple
least singular value does not interrupt the continuation.  The folded pair
\(u=2z^2-1\), \(v=2y^2-1\) and its reconstructed half variables are
continuous, while the sheet inequalities are strict.  Thus \(\mathcal C\)
is relatively open in the gap.

It is also relatively closed.  At a relative closure point the chord
identity and all sheet inequalities persist weakly.  The possible closure
cases are as follows.  (1) If \(\Delta_*(x)=0\), then the reconstructed half
variables coincide, which is impossible because
\[
 z\geq u_K^{\max}>\cos\frac{\pi}{K}
 \geq\cos\frac{j\pi}{K}\geq y.
\]
(2) Neither half variable can change sign in a noncentral positive gap:
\(z>0\), while \(y=0\) would force \(x=0\), outside the gap.  (3) The outer
variable cannot reach \(\mathfrak c_h\) in the interior, because
\eqref{eq:channel-xs} would then give \(s=0\).  (4) The inner variable cannot
reach a nodal endpoint: the common level would be zero, the endpoint-side
solution would be \(\mathfrak c_h\), and \eqref{eq:channel-xs} would give
\(s=0\), whereas the middle root has no zero inside a spectral gap.  (5) The
outer variable cannot leave its endpoint-side component, since
its first such point would be \(u_K^{\max}\), contradicting
\eqref{eq:outer-level-dominates}.  Passage through \(z=1\) is harmless; only
its coordinate changes.  Hence every relative closure point remains in
\(\mathcal C\).  Since the open gap is an interval, \(\mathcal C\) is the
entire gap.  This proves the first exhaustion direction: every actual gap
point lies on the selected endpoint-side chord sheet.

For the converse direction, the reconstructed angle \(\Theta\) extends
continuously to the closed spectral gap.  The exact folded variables at the
two endpoints give
\[
 \Theta(\lambda_{K-1,j+1})=\frac{(j+1)\pi}{K},\qquad
 \Theta(\lambda_{K-1,j})=\frac{j\pi}{K}.
\]
For every \(\theta_0\in(j\pi/K,(j+1)\pi/K)\), the intermediate value
theorem therefore supplies an interior gap point \(x\) with
\(\Theta(x)=\theta_0\).  The classification just proved and uniqueness of
the endpoint-side outer solution then give
\[
 y(x)=\cos\theta_0,\qquad
 z(x)=z_K(\mathcal W_K(\cos\theta_0)),
\]
so the corresponding abscissa and signed root are exactly those of the
prescribed $j$-chord.  Thus every chord parameter produces an actual point of
the ordered middle branch.  By \cref{lem:middle-branch}, its positive
magnitude is \(s_{\min}(xI_{K-1}-A_{K-1})\).  Both exhaustion directions
have now been proved without assuming that \(\theta\mapsto x\) is monotone.

When \(K\) is odd, the same argument applies to the positive half of
the single central gap, continuing from the positive spectral endpoint
to \(x=0\).  Because \(z>0\), the identity
\(x=(2r/\mathfrak c_h)\cos\theta\,z\) forces \(\theta=\pi/2\) at
\(x=0\).
Thus every parameter on that half occurs, and reflection covers the
negative half.

\medskip\noindent\emph{Noncentral gap comparison.}\quad
We now compare the noncentral middle branches.  Let \(N\geq3\), put
\(M=N+1\), and set \(\kappa=N/M\).  Every positive noncentral gap of
the \(M\)-problem
satisfies
\begin{equation}\label{eq:noncentral-domain}
 j\geq1,\qquad \frac{(j+1)\pi}{N}\leq\frac{\pi}{2}.
\end{equation}
Indeed, such gaps satisfy \(j+1\leq N/2\).  Hence
\(\theta_0=\theta/\kappa\) belongs to the
same-indexed \(N\)-nodal interval, and every trigonometric argument
below lies in \((0,\pi/2]\).  Consider a point of such a branch for the
matrix of order \(M-1\), with
\[
 \theta\in\left(\frac{j\pi}{M},\frac{(j+1)\pi}{M}\right),
 \qquad \frac{(j+1)\pi}{N}\leq\frac{\pi}{2}.
\]
First suppose \(\Phi_M(\theta)=\Psi_M(\eta)\).  Put
\(\theta_0=\theta/\kappa\).  Then \(\theta_0\) belongs to the
corresponding \(N\)-nodal interval.  Concavity of sine gives
\begin{equation}\label{eq:Phi-scale}
 \frac{\Phi_N(\theta_0)}{\Phi_M(\theta)}
 =\frac{\mathcal F_h(\theta_0)}{\mathcal F_h(\theta)}>\kappa.
\end{equation}
Indeed, \(\theta=\kappa\theta_0\) and strict concavity gives
\(\sin\theta>\kappa\sin\theta_0\); hence
\[
 \mathcal F_h(\theta_0)^2-\kappa^2\mathcal F_h(\theta)^2
 =(1-\kappa^2)+\mathfrak s_h^2\left(
 \frac1{\sin^2\theta_0}-\frac{\kappa^2}{\sin^2\theta}
 \right)>0.
\]
On the other hand,
\begin{equation}\label{eq:Psi-scale}
 \frac{\Psi_N(\eta)}{\Psi_M(\eta)}
 =\frac{\sinh(N\eta)}{\sinh(M\eta)}\leq\kappa,
\end{equation}
with strict inequality for \(\eta>0\), by convexity of \(\sinh\) on
the positive half-axis.  Hence
\(\Phi_N(\theta_0)>\Psi_N(\eta)\).  Since \(\Phi_N\) vanishes at the right
endpoint \((j+1)\pi/N\), its rightmost solution
\(\widehat\theta\) of
\[
 \Phi_N(\widehat\theta)=\Psi_N(\eta)
\]
exists by the intermediate value theorem and satisfies
\(\widehat\theta>\theta_0>\theta\).  Strict log-concavity of
\(\mathcal W_N\) makes it the unique solution on the descending half
of that lobe.  The same exterior
variable is endpoint-side for the \(N\)-problem as well.  Indeed, for
\(u=\cosh\eta>1\),
\[
 \frac{d}{du}\log\mathcal W_M(u)-
 \frac{d}{du}\log\mathcal W_N(u)
 =\frac{M\coth(M\eta)-N\coth(N\eta)}{\sinh\eta}>0.
\]
The numerator is positive because, for \(\eta>0\),
\[
 \frac{d}{dt}\bigl(t\coth(t\eta)\bigr)
 =\frac{\sinh(t\eta)\cosh(t\eta)-t\eta}{\sinh^2(t\eta)}>0.
\]
The endpoint-side branch is characterized by a nonpositive logarithmic
derivative.  The displayed difference therefore makes the
\(N\)-derivative strictly negative.  The case \(\eta=0\) follows by
continuity.
The chord classification above says that the constructed value
\(\mathfrak s_{\rm hyp}(\widehat\theta,\eta)\) is the least singular
value at its constructed abscissa.  Formula
\eqref{eq:channel-xs} and the strict increase of sine on the interval
now give
\begin{equation}\label{eq:hyper-height-increase}
 \mathfrak s_{\rm hyp}(\widehat\theta,\eta)
 >\mathfrak s_{\rm hyp}(\theta,\eta).
\end{equation}

Suppose next that the exterior endpoint is elliptic, so that
\(\Phi_M(\theta)=\Phi_M(\phi)\), \(0<\phi<\pi/M\).  (The case
\(\phi=0\)
is the limiting hyperbolic case already covered.)  Put
\[
 \theta_0=\theta/\kappa,\qquad \phi_0=\phi/\kappa.
\]
The ratio
\[
 \mathcal A_{N,M}^{\Phi}(t)
 =\frac{\Phi_N(t/\kappa)}{\Phi_M(t)}
 =\frac{\mathcal F_h(t/\kappa)}{\mathcal F_h(t)}
\]
is strictly increasing while \(0<t<t/\kappa\leq\pi/2\).  Indeed,
\[
 \frac{d}{dt}\log\mathcal A_{N,M}^{\Phi}(t)
 =\mathfrak d_h(t)-\kappa^{-1}\mathfrak d_h(t/\kappa)>0,\qquad
 \mathfrak d_h(t)=\frac{\mathfrak s_h^2\cot t}
 {\mathfrak s_h^2+\sin^2t},
\]
because
\[
 t\mathfrak d_h(t)=(t\cot t)
 \frac{\mathfrak s_h^2}{\mathfrak s_h^2+\sin^2t}
\]
is strictly decreasing: both positive factors decrease on
\((0,\pi/2)\), and the second decreases strictly.  Since
\(\theta>\pi/M>\phi\), it follows that
\(\Phi_N(\theta_0)>\Phi_N(\phi_0)\).  The rightmost solution of
\[
 \Phi_N(\widehat\theta)=\Phi_N(\phi_0)
\]
therefore has \(\widehat\theta>\theta_0\).  Its exterior variable is
endpoint-side.  Indeed, if
\(\Lambda_K(t)=(\log\Phi_K)'(t)\), then on the
endpoint-side part of the first outer lobe
\(\Lambda_M(\phi)\geq0\), and differentiation of
\(\Phi_N(t/\kappa)=\mathcal A_{N,M}^{\Phi}(t)\Phi_M(t)\) gives
\[
 \kappa^{-1}\Lambda_N(t/\kappa)
 =\Lambda_M(t)
 +\frac{(\mathcal A_{N,M}^{\Phi})'(t)}{\mathcal A_{N,M}^{\Phi}(t)}>0.
\]
Because
\(\Lambda_K(\phi)=-\sin\phi\,(\log\mathcal W_K)'(\cos\phi)\),
positivity of \(\Lambda_N(\phi_0)\) is exactly the endpoint-side
condition for the
\(N\)-outer solution.  The one-maximum property now fixes the label, and the chord
classification says again that
\(\mathfrak s_{\rm ell}(\widehat\theta,\phi_0)\) is the least singular
value at the constructed abscissa.  Here
\(\widehat\theta>\theta_0>\theta\) and \(\phi_0>\phi\), with all four
angles below \(\pi/2\).  Thus both sine factors in the elliptic version
of \eqref{eq:channel-xs} have increased, at least one strictly, and hence
\begin{equation}\label{eq:elliptic-height-increase}
 \mathfrak s_{\rm ell}(\widehat\theta,\phi_0)
 >\mathfrak s_{\rm ell}(\theta,\phi).
\end{equation}

\medskip\noindent\emph{Central-gap interlacing.}\quad
It remains to compare the central sheets.  We do this in the next
calculation.

For \(m\geq1\), let
\begin{equation}\label{eq:c-height-def}
 c_m=s_{\min}(aI_m+V_m).
\end{equation}
By reflection, the two gaps adjacent to the zero eigenvalue of
\(A_{2m+1}\) have the same height; denote it by
\begin{equation}\label{eq:d-height-def}
 d_m=\max_{0\leq x\leq2r\sin\frac{\pi}{2m+2}}
 s_{\min}(xI_{2m+1}-A_{2m+1}).
\end{equation}
We claim
\begin{equation}\label{eq:central-interlace}
 c_m>d_m>c_{m+1}.
\end{equation}
Put
\[
 \zeta=\frac{a+a^{-1}}2.
\]
For every \(j\geq1\), there is a unique
\(\rho_j>\cos\frac{\pi}{j+1}\) such that
\begin{equation}\label{eq:rho-central}
 aU_j(\rho_j)=U_{j-1}(\rho_j).
\end{equation}
Indeed, the Christoffel--Darboux identity gives, above the largest zero
of \(U_j\),
\[
 \frac{d}{dt}\frac{U_{j-1}(t)}{U_j(t)}
 =-\frac{2\sum_{k=0}^{j-1}U_k(t)^2}{U_j(t)^2}<0.
\]
The ratio decreases from \(+\infty\) to \(0\), proving existence and
uniqueness.  A three-term determinant recurrence gives
\begin{equation}\label{eq:c-determinant}
 \det\left((aI_j+V_j)^*(aI_j+V_j)
 -(1+a^2-2at)I_j\right)
 =a^j\left(U_j(t)-a^{-1}U_{j-1}(t)\right).
\end{equation}
Thus the largest root \(\rho_j\) gives the least eigenvalue.  Together
with
\[
 U_j(t)^2+U_{j-1}(t)^2-2tU_j(t)U_{j-1}(t)=1,
\]
equations \eqref{eq:rho-central} and \eqref{eq:c-determinant} give
\begin{equation}\label{eq:c-rho-identities}
 c_j^2=1+a^2-2a\rho_j=\frac1{U_j(\rho_j)^2},\qquad
 \zeta-\rho_j=\frac1{2U_{j-1}(\rho_j)U_j(\rho_j)}.
\end{equation}

We shall also use that \(c_m\) is exactly the height of the central gap
of \(A_{2m}\).  Let
\[
 P_{\rm oe}=[\mathbf e_1,\mathbf e_3,\ldots,\mathbf e_{2m-1},
 \mathbf e_2,\mathbf e_4,\ldots,\mathbf e_{2m}]
\]
be the permutation matrix that places odd coordinates first.  Then
\[
 P_{\rm oe}^T A_{2m}P_{\rm oe}=
 \begin{bmatrix}0&I_m+aV_m^T\\aI_m+V_m&0\end{bmatrix}.
\]
Moreover,
\[
 (I_m+aV_m^T)^*(I_m+aV_m^T)
 -J_m(aI_m+V_m)^*(aI_m+V_m)J_m
 =(1-a^2)\mathbf e_m\mathbf e_m^T\geq0.
\]
The singular values of the block off-diagonal matrix are the union of
the singular values of its two off-diagonal blocks.  The Gram
inequality shows that
\(s_{\min}(I_m+aV_m^T)\geq s_{\min}(aI_m+V_m)=c_m\), and hence
\(s_{\min}(A_{2m})=c_m\).  To prove that this is the
maximum throughout the central gap, put \(K=2m+1\) and use the half
variables
\[
 y=\cos\theta,\qquad
 z=\cosh\eta\quad\hbox{or}\quad z=\cos\phi.
\]
On the positive half of the central middle branch, the preceding chord
classification gives
\begin{equation}\label{eq:even-central-half}
 \mathcal W_K(y)=\mathcal W_K(z),\qquad
 x=\frac{2r}{\mathfrak c_h}yz,\qquad
 s^2=\frac{4a}{\mathfrak c_h^2}
 (\mathfrak c_h^2-y^2)(\mathfrak c_h^2-z^2),
\end{equation}
where \(0\leq y<\cos\frac{m\pi}{K}\) is in the central lobe and
\(z\) is the
endpoint-side outer solution.  The function \(\mathcal W_K\) is even,
and strict log-concavity on the central lobe gives
\[
 \mathcal W_K(y)<\mathcal W_K(0)=\mathfrak c_h\qquad(y>0).
\]
Let
\[
 z_0=\sqrt{\frac{1+\rho_m}{2}}.
\]
The half-angle identity and \eqref{eq:rho-central} give
\[
 U_{2m}(z_0)=U_m(\rho_m)+U_{m-1}(\rho_m)
 =(1+a)U_m(\rho_m).
\]
Together with \eqref{eq:c-rho-identities}, this yields
\begin{equation}\label{eq:z0-central}
 \mathcal W_K(z_0)=\mathfrak c_h,\qquad
 4a(\mathfrak c_h^2-z_0^2)=2a(\zeta-\rho_m)=c_m^2.
\end{equation}
At \(x=0\), \cref{lem:middle-branch} and the chord-exhaustion argument
identify the nearest root with \(c_m\), so \(z_0\) is the endpoint-side
outer solution of level \(\mathfrak c_h\).  If \(x\neq0\), the level
in \eqref{eq:even-central-half} is strictly below \(\mathfrak c_h\), and hence its
endpoint-side solution satisfies \(z>z_0\).  Equations
\eqref{eq:even-central-half}--\eqref{eq:z0-central} then give
\[
 s^2<4a(\mathfrak c_h^2-z_0^2)=c_m^2.
\]
Reflection handles the negative half of the gap.  Thus its height is
exactly \(c_m\), attained only at \(x=0\).

\medskip\noindent\emph{The upper comparison \(d_m<c_m\).}\quad
Set
\[
 \omega_m^{\rm e}=\cos\frac{\pi}{m+1},\qquad
 \mathcal V_m(t)=\sqrt{(\zeta-t)(1+t)}\,U_m(t).
\]
Here are the details of the central specialization.  Take \(K=2m+2\)
and set
\[
 t_{\rm in}=2y^2-1=\cos(2\theta),\qquad
 t_{\rm out}=2z^2-1.
\]
On the positive central gap,
\(-1<t_{\rm in}<-\omega_m^{\rm e}\), while the endpoint-side outer variable
satisfies
\(\omega_m^{\rm e}<t_{\rm out}<\zeta\), where
\(\zeta=2\mathfrak c_h^2-1\).  The identity
\(U_{2m+1}(t)=2tU_m(2t^2-1)\) gives
\[
 \mathcal W_{2m+2}(y)^2
  = (\zeta-t_{\rm in})(1+t_{\rm in})U_m(t_{\rm in})^2,
 \qquad
 \mathcal W_{2m+2}(z)^2
  = (\zeta-t_{\rm out})(1+t_{\rm out})U_m(t_{\rm out})^2.
\]
Because \(\zeta+1=2\mathfrak c_h^2\), the formulas
\(x=(2r/\mathfrak c_h)yz\) and
\(s^2=(4a/\mathfrak c_h^2)
(\mathfrak c_h^2-y^2)(\mathfrak c_h^2-z^2)\) become the following chord
parametrization:
\begin{equation}\label{eq:central-chord}
\begin{gathered}
 -1<t_{\rm in}<-\omega_m^{\rm e},\qquad
 \omega_m^{\rm e}<t_{\rm out}<\zeta,\qquad
 \mathcal V_m(t_{\rm out})=|\mathcal V_m(t_{\rm in})|,\\
 x^2=\frac{2a(1+t_{\rm out})(1+t_{\rm in})}{\zeta+1},\qquad
 s^2=\frac{2a(\zeta-t_{\rm out})(\zeta-t_{\rm in})}{\zeta+1},
\end{gathered}
\end{equation}
where \(t_{\rm out}\) is the descending, or largest, solution.  To see this
from \eqref{eq:chord-equation}, use
\[
 \mathfrak s_h^2+\sin^2\theta=\frac{\zeta-t_{\rm in}}{2},\qquad
 \mathfrak s_h^2-\sinh^2\eta=\frac{\zeta-t_{\rm out}}{2}.
\]
Continuation through \(t_{\rm out}=1\) uses the elliptic coordinate
\(z=\cos\phi\).

Let \(\rho=\rho_m\).  The function \(\log\mathcal V_m\) is strictly
concave on \((\omega_m^{\rm e},\zeta)\).  If
\((\log\mathcal V_m)'(\rho)\geq0\), the descending root
satisfies \(t_{\rm out}\geq\rho\), because \(\rho\) lies at or before the unique
lobe maximum and the descending root lies at or after it.  Suppose instead that
\((\log\mathcal V_m)'(\rho)<0\).  Then \(\rho>a\).  Indeed, otherwise
\(\rho\leq a<1\), so \(\rho<1\) and in particular \(\rho^2-1\neq0\).  The
identity
\[
 \frac{U_m'(\rho)}{U_m(\rho)}
 =\frac{m\rho-(m+1)a}{\rho^2-1}
\]
would give
\[
\begin{aligned}
 (\log\mathcal V_m)'(\rho)
 &=\frac{(2m+2)a+1-(2m+1)\rho}{2(1-\rho^2)}
   -\frac1{2(\zeta-\rho)}\\
 &\geq\frac1{2(1-\rho)}-\frac1{2(\zeta-\rho)}\geq0,
\end{aligned}
\]
a contradiction.  Thus the displayed quotient has been used only on its
valid domain; if \(\rho=1\), the desired conclusion \(\rho>a\) is immediate
from \(a<1\) and no quotient is invoked.
Writing \(t_{\rm in}=-\cos\phi_{\rm c}\),
\(0<\phi_{\rm c}<\pi/(m+1)\), gives
\[
\begin{aligned}
 |\mathcal V_m(t_{\rm in})|^2
 &=\frac{(\zeta+\cos\phi_{\rm c})
          \sin^2((m+1)\phi_{\rm c})}{1+\cos\phi_{\rm c}}\\
 &<\frac{\zeta+\omega_m^{\rm e}}{1+\omega_m^{\rm e}}
 <\frac{1+a}{2a}<\frac{1+\rho}{2a}=\mathcal V_m(\rho)^2.
\end{aligned}
\]
Now \(\rho\) is on the descending side and
\(|\mathcal V_m(t_{\rm in})|<\mathcal V_m(\rho)\), so the
unique descending root again satisfies \(t_{\rm out}>\rho\).  In either case,
\(t_{\rm out}\geq\rho\), and \(t_{\rm in}>-1\) makes the following
inequality strict:
\[
 s^2<2a(\zeta-\rho)=c_m^2.
\]
This holds along the whole central sheet, so \(d_m<c_m\).

\medskip\noindent\emph{The lower comparison \(d_m>c_{m+1}\).}\quad
Put \(L=m+1\geq2\), and abbreviate
\[
 \rho=\rho_L,\quad u_-=U_{L-1}(\rho),\quad u_+=U_L(\rho),\quad
 \omega_L=\cos\frac{\pi}{2L},\quad \Delta_\omega=1-\omega_L,
\]
\begin{equation}\label{eq:central-D-xi}
 \mathfrak D_L=1+2(\rho+\omega_L)u_-u_+,\qquad
 \xi=\rho-\frac{\Delta_\omega}{\mathfrak D_L}.
\end{equation}
Thus \(a=u_-/u_+\) and \(c_L=1/u_+\).

\medskip\noindent\emph{The decisive scalar estimate.}\quad
We first establish
\begin{equation}\label{eq:central-target}
 (1+\omega_L)(1+\xi)^2U_{L-1}(\xi)^2
 <(\rho+\omega_L)\mathfrak D_L.
\end{equation}

Put \(b=u_+/u_-=1/a\), and define
\[
 Q_b=b^2-2b\rho+1,\qquad S_b=b^2+2b\omega_L+1.
\]
The Chebyshev identity gives \(u_-^2Q_b=1\), and
\begin{equation}\label{eq:central-D-QS}
 \mathfrak D_L=S_b/Q_b,\qquad
 \Delta_\xi:=\rho-\xi=\Delta_\omega Q_b/S_b.
\end{equation}
Indeed, \(u_+=bu_-\) and \(u_-^2=Q_b^{-1}\), so
\[
 \mathfrak D_L=1+2(\rho+\omega_L)u_-u_+
  =1+\frac{2b(\rho+\omega_L)}{Q_b}=\frac{S_b}{Q_b},
\]
and the formula for \(\Delta_\xi\) follows from
\eqref{eq:central-D-xi}.
The ratio \(U_L/U_{L-1}\) is strictly increasing above the largest zero
of \(U_{L-1}\), and equals \(1\) at
\(\rho_{\rm ref}=\cos\frac{\pi}{2L+1}\).  Hence
\(\rho>\rho_{\rm ref}\).  Since
\(S_b-Q_b=2b(\rho+\omega_L)>0\),
\[
 \xi>\rho_{\rm ref}-\Delta_\omega>2\omega_L-1
 >2\omega_L^2-1=\cos\frac{\pi}{L}.
\]
Therefore \(f_L(x)=(1+x)U_{L-1}(x)\) is positive, strictly increasing,
and strictly log-concave on \([\xi,\rho]\).  In particular,
\(\mathfrak d_\rho=(\log f_L)'(\rho)>0\), and
\[
 \mathfrak d_\rho=
 \begin{cases}
  \dfrac{L(b-\rho)-1}{\rho^2-1},&\rho\neq1,\\[6pt]
  \dfrac{2L^2+1}{6},&\rho=1.
 \end{cases}
\]
The endpoint value follows directly from \(U_{L-1}(1)=L\) and
\(U_{L-1}'(1)=L(L^2-1)/3\).  Moreover,
\[
 \log f_L(\rho)-\log f_L(\xi)>\Delta_\xi\mathfrak d_\rho,
 \qquad
 f_L(\xi)^2<f_L(\rho)^2e^{-2\Delta_\xi\mathfrak d_\rho}
 <\frac{f_L(\rho)^2}{1+2\Delta_\xi\mathfrak d_\rho}.
\]
Here the first strict inequality is the integral form of strict
log-concavity.  Since
\(f_L(\rho)^2=(1+\rho)^2/Q_b\) and
\(\Delta_\xi=\Delta_\omega Q_b/S_b\), substitution into
\eqref{eq:central-target} reduces it to
\begin{equation}\label{eq:N-positive-target}
 \mathcal N=(\rho+\omega_L)
 (S_b+2\Delta_\omega Q_b\mathfrak d_\rho)
 -(1+\omega_L)(1+\rho)^2>0.
\end{equation}
We shall repeatedly use
\begin{equation}\label{eq:Delta-omega-basic-bound}
 \Delta_\omega=1-\cos\frac{\pi}{2L}
 <\frac{\pi^2}{8L^2}<\frac5{4L^2},
 \qquad \Delta_\omega L<1.
\end{equation}

Suppose first that \(\rho>1\).  Put \(u=b-\rho\) and
\(\Delta_+=\rho^2-1\).  Since \(\mathfrak d_\rho>0\), \(u>1/L\).
Keeping \(\rho\) fixed,
direct differentiation gives
\[
 \frac{\partial_b\mathcal N}{2(\rho+\omega_L)}
 =\rho+\omega_L+u-\Delta_\omega L
 +\frac{\Delta_\omega(3Lu^2-2u)}{\Delta_+}.
\]
This is strictly increasing in \(u\geq1/L\), and its value at \(1/L\)
is positive because \(\Delta_\omega L<1\).  Thus \(\mathcal N\)
exceeds its value at
\(b=\rho+1/L\).  With
\(s_0=\rho+\omega_L\) and \(\omega_+=1+\omega_L\), that value is
\[
 \Upsilon(s_0)
 =s_0\left[\left(s_0+\frac1L\right)^2+\omega_+\Delta_\omega\right]
 -\omega_+(s_0+\Delta_\omega)^2.
\]
For \(s_0\geq\omega_+\),
\[
 \Upsilon''(s_0)=6s_0+4/L-2\omega_+>0,\qquad
 \Upsilon'(\omega_+)=2\omega_L\omega_++4\omega_+/L+1/L^2>0,
\]
and
\[
 \Upsilon(\omega_+)
 =\omega_+(-2\Delta_\omega+2\omega_+/L+1/L^2)>0.
\]
Thus \eqref{eq:N-positive-target} holds when \(\rho>1\).  If \(\rho=1\),
then \(b=(L+1)/L\), \(Q_b=L^{-2}\), and direct substitution, using the
endpoint value of \(\mathfrak d_\rho\), gives
\[
 \mathcal N=\Upsilon(1+\omega_L)
 +\frac{2(1+\omega_L)\Delta_\omega}{L^2}
   \frac{2L^2+1}{6}>0.
\]
The first term is positive by the preceding calculation, and every factor
in the correction term is positive.

Suppose next that \(\rho<1\).  Set
\[
 \upsilon=1-\rho,\qquad \Delta_b=b-1,\qquad
 u=\Delta_b+\upsilon,\qquad \Delta_-=\upsilon(2-\upsilon).
\]
Writing \(\rho=\cos\theta_\rho\) gives
\(0<L\theta_\rho<\pi/2\) and
\[
 u=b-\rho
 =\frac{\sin((L+1)\theta_\rho)}{\sin(L\theta_\rho)}
  -\cos\theta_\rho
 =\sin\theta_\rho\cot(L\theta_\rho).
\]
Because \(b>1\), \(u=\Delta_b+\upsilon>\upsilon>0\).  Also
\(\tan(L\theta_\rho)>L\tan\theta_\rho>L\sin\theta_\rho\), since
\(x\mapsto\tan x/x\) is strictly increasing on \((0,\pi/2)\).
Hence
\begin{equation}\label{eq:u-small-range}
 0<\upsilon<u<1/L.
\end{equation}
With \(\rho\) fixed, write \(\mathcal N_\rho^{<}(\mathord\cdot)\) for
\(\mathcal N\) viewed as a function of \(\Delta_b=b-1\).
A direct expansion of \eqref{eq:N-positive-target} gives
\[
 \mathcal N_\rho^{<}(\Delta_b)-\mathcal N_\rho^{<}(0)
 =\frac{\Delta_b(\rho+\omega_L)}{\Delta_-}\mathfrak q(u),
\]
where
\[
 \mathfrak q(u)=\Delta_-[2(2-\Delta_\omega)+u-\upsilon]
 +2\Delta_\omega
 [u+\upsilon-L(u^2+u\upsilon+\upsilon^2)-L\Delta_-].
\]
This is a concave quadratic.  Moreover,
\[
 \mathcal N_\rho^{<}(0)=\frac{\upsilon}{2-\upsilon}
 \bigl[(2-\Delta_\omega)(2-\upsilon)^2
 -2\Delta_\omega^2
 -4\Delta_\omega L(2-\upsilon-\Delta_\omega)\bigr].
\]
The following elementary bounds hold; the second follows from
\(\rho>\rho_{\rm ref}\) and
\(1-\rho_{\rm ref}<\pi^2/[2(2L+1)^2]<5/(2L+1)^2\):
\begin{equation}\label{eq:Delta-upsilon-bounds}
 \Delta_\omega<\frac5{4L^2},\qquad
 \upsilon<\frac5{(2L+1)^2}\leq\frac15.
\end{equation}
The last inequality is an equality when \(L=2\); because the preceding
bound for \(\upsilon\) is strict, one still has \(\upsilon<1/5\).  These bounds
show that the bracketed expression is positive, since it equals
\[
 (2-\upsilon)
 \bigl[(2-\Delta_\omega)(2-\upsilon)-4\Delta_\omega L\bigr]
 +2\Delta_\omega^2(2L-1)
\]
and, using \(L\geq2\),
\[
 (2-\Delta_\omega)(2-\upsilon)
 >\frac{27}{16}\frac95=\frac{243}{80}
 >\frac52\geq4\Delta_\omega L.
\]
At the endpoints of
\eqref{eq:u-small-range},
\[
 \mathfrak q(\upsilon)=\upsilon
 \bigl[(2-\upsilon)(4-2\Delta_\omega)
 +4\Delta_\omega-4\Delta_\omega L(1+\upsilon)\bigr]>0,
\]
\[
 \mathfrak q(1/L)=\upsilon
 \bigl[(2-\upsilon)(4-2\Delta_\omega+1/L-\upsilon)
 -4\Delta_\omega L\bigr]>0.
\]
Indeed, in these two lines the positive leading terms exceed
\[
 \frac95\left(4-\frac58\right)=\frac{243}{40},
 \qquad
 \frac95\left(4-\frac58-\frac15\right)=\frac{1143}{200},
\]
while the negative terms have magnitudes less than
\((5/L)(6/5)\leq3\) and \(5/L\leq5/2\), respectively.  Concavity now gives
\(\mathfrak q(u)>0\), proving \eqref{eq:N-positive-target}, and hence
\eqref{eq:central-target}, in all cases.

\medskip\noindent\emph{Certification of the least singular value.}\quad
Since \(a>0\), \(\Delta_\omega>0\), \(\rho+\omega_L>0\), and
\(\zeta+\omega_L>0\), let \(x_*>0\) be the positive square root determined
by
\begin{equation}\label{eq:x-star-central}
 x_*^2=\frac{2a\Delta_\omega(\rho+\omega_L)}{\zeta+\omega_L},
 \qquad c=c_L.
\end{equation}
At \(x=x_*\) and \(s=\pm c\), the real pair
\(u=\xi\), \(v=-\omega_L\) satisfies the two symmetric relations in
\eqref{eq:fold-characteristic}.  Indeed, its sum and product are
\(\xi-\omega_L\) and \(-\omega_L\xi\), and the required identities
follow from
\[
 \mathfrak D_L=\frac{\zeta+\omega_L}{\zeta-\rho},\qquad
 \xi=\frac{\rho(\zeta+1)+\zeta(\omega_L-1)}{\zeta+\omega_L}.
\]
Moreover, \(\xi+\omega_L>0\), so the distinct branch of the divided
difference in \eqref{eq:fold-odd} applies.
For convenience, write the folded numerator as
\[
 \mathfrak n_s(t)=-[x(\zeta-t)+s(1+t)]U_{L-1}(t).
\]
Put
\[
 \mathsf{X}_\xi=x_*(\zeta-\xi),\quad
 \mathsf{Y}_\xi=c(1+\xi),\quad
 \mathsf{X}_\omega=x_*(\zeta+\omega_L),\quad
 \mathsf{Y}_\omega=c\Delta_\omega,
\]
and put
\(U_\xi=U_{L-1}(\xi)\),
\(U_{-\omega}=U_{L-1}(-\omega_L)\).  The identities
\[
 \mathfrak D_L=2a(\zeta+\omega_L)u_+^2,\qquad
 \mathfrak D_L(\rho+\omega_L)(\zeta-\xi)
 =(\zeta+\omega_L)(1+\xi)
\]
give
\(\mathsf{Y}_\omega\mathsf{Y}_\xi
=\mathsf{X}_\omega\mathsf{X}_\xi\).  The identity \(u_-^2Q_b=1\) gives
\(Q_b>0\), and \(S_b-Q_b=2b(\rho+\omega_L)>0\), so
\(\mathfrak D_L=S_b/Q_b>1\).  Moreover, \(L\geq2\) gives
\(\omega_L>1/2\), while \(\rho>\rho_{\rm ref}>\omega_L\); hence
\[
 \rho+\omega_L>2\omega_L>1>1-\omega_L>0.
\]
Consequently,
\[
 \frac{\mathsf{X}_\omega^2}{\mathsf{Y}_\omega^2}
 =\frac{\mathfrak D_L(\rho+\omega_L)}{1-\omega_L}>1,
\]
so, with
\(\mathfrak r_*=\mathsf{X}_\xi/\mathsf{Y}_\xi
=\mathsf{Y}_\omega/\mathsf{X}_\omega\in(0,1)\), the common nonzero factor
in \eqref{eq:fold-odd} is \(a^{L-1}/(\xi+\omega_L)\), and the two numerator
differences are
\[
\begin{aligned}
 \mathfrak n_c(\xi)-\mathfrak n_c(-\omega_L)
  &=(1+\mathfrak r_*)
    (\mathsf{X}_\omega U_{-\omega}-\mathsf{Y}_\xi U_\xi),\\
 \mathfrak n_{-c}(\xi)-\mathfrak n_{-c}(-\omega_L)
  &=(1-\mathfrak r_*)
    (\mathsf{X}_\omega U_{-\omega}+\mathsf{Y}_\xi U_\xi).
\end{aligned}
\]
Their product is positive by \eqref{eq:central-target}, because
\[
 U_{-\omega}^2=\frac1{1-\omega_L^2},\qquad
 \mathsf{X}_\omega^2U_{-\omega}^2
 =\frac{c^2\mathfrak D_L(\rho+\omega_L)}{1+\omega_L}
 >\mathsf{Y}_\xi^2U_\xi^2.
\]
Thus
\begin{equation}\label{eq:q-product-central}
 q_{2L-1}(x_*,c)q_{2L-1}(x_*,-c)>0.
\end{equation}

At \(x=0\), the second-smallest singular value of \(A_{2L-1}\) is
\[
 s_{L,2}^{(0)}=\sqrt{1+a^2-2a\cos\frac{\pi}{L}}.
\]
Indeed, after separating odd and even coordinates, the two rectangular
off-diagonal blocks have Gram matrices equal to the
\((L-1)\)-dimensional tridiagonal Toeplitz matrix with diagonal
\(1+a^2\) and off-diagonal \(a\).  Their singular values therefore have
squares
\[
 1+a^2+2a\cos\frac{k\pi}{L},\qquad 1\leq k\leq L-1,
\]
and each occurs twice, while the remaining singular value is zero.
We claim
\begin{equation}\label{eq:x-c-second-singular}
 x_*+c<2\sqrt a\sin\frac{\pi}{2L}\leq s_{L,2}^{(0)}.
\end{equation}
Indeed, recall \(\rho_{\rm ref}=\cos\frac{\pi}{2L+1}\).  Since
\(\rho>\rho_{\rm ref}\) and \(U_{L-1}\) is strictly increasing above its
largest zero,
\[
 u_-=U_{L-1}(\rho)>U_{L-1}(\rho_{\rm ref})=:u_{-,\rm ref},
 \qquad u_{-,\rm ref}^{-2}=2(1-\rho_{\rm ref}).
\]
The identities \(c=a/u_-\) and \(c^2=2a(\zeta-\rho)\) imply
\(\zeta-\rho<a(1-\rho_{\rm ref})<1-\rho_{\rm ref}\).  Also
\[
 \frac{x_*^2}{2a}
 =\Delta_\omega\frac{\rho+\omega_L}{\zeta+\omega_L}
 <\Delta_\omega.
\]
Put
\[
 \theta_L=\frac{\pi}{4L},\qquad
 \Delta_L=\frac{\theta_L}{2L+1}.
\]
Then \(\Delta_\omega=2\sin^2\theta_L\) and
\(1-\rho_{\rm ref}=2\sin^2(\theta_L-\Delta_L)\).  Since all quantities are
positive, the preceding quotient bound for \(x_*^2\) gives
\[
 x_*<2\sqrt a\sin\theta_L.
\]
Likewise,
\[
 c^2=2a(\zeta-\rho)
 <2a^2(1-\rho_{\rm ref})
 =4a^2\sin^2(\theta_L-\Delta_L),
\]
and hence
\[
 c<2a\sin(\theta_L-\Delta_L)
 <2\sqrt a\sin(\theta_L-\Delta_L),
\]
where the last inequality uses \(0<a<1\), hence \(a<\sqrt a\).
It remains to prove
\begin{equation}\label{eq:angle-inequality}
 \sin\theta_L+\sin(\theta_L-\Delta_L)<\sin(2\theta_L).
\end{equation}
Now
\[
 2\sin\theta_L-\sin(2\theta_L)
 =2\sin\theta_L(1-\cos\theta_L)<\theta_L^3
\]
and
\[
 \sin\theta_L-\sin(\theta_L-\Delta_L)
 =\int_{\theta_L-\Delta_L}^{\theta_L}\cos u\,du
 >\Delta_L\cos\theta_L.
\]
Thus \eqref{eq:angle-inequality} follows from
\(\Delta_L\cos\theta_L>\theta_L^3\), or
\(2\cos\theta_L>\theta_L(\pi+2\theta_L)\).  For
\(0<\theta_L\leq\pi/8\), using \(\pi^2<10\),
\[
 2\cos\theta_L>2-\theta_L^2\geq\frac{59}{32}
 >\frac{25}{16}\geq\theta_L(\pi+2\theta_L).
\]
Consequently \(x_*+c<2\sqrt a\sin\frac{\pi}{2L}\).  Finally,
\[
 (s_{L,2}^{(0)})^2-4a\sin^2\frac{\pi}{2L}
 =(1-a)^2\geq0,
\]
which proves \eqref{eq:x-c-second-singular}.

Let \(\mu_1,\ldots,\mu_{2L-1}\) be the eigenvalues of the real symmetric
matrix
\[
 B_{2L-1}(x_*)=(x_*I_{2L-1}-A_{2L-1})J_{2L-1}.
\]
Right multiplication by \(J_{2L-1}\) preserves singular values, so the
numbers \(|\mu_i|\) are precisely the singular values of
\(x_*I_{2L-1}-A_{2L-1}\).  Moreover,
\[
 q_{2L-1}(x_*,c)q_{2L-1}(x_*,-c)
 =\prod_{i=1}^{2L-1}(\mu_i^2-c^2).
\]
The positive product in \eqref{eq:q-product-central} is nonzero and
therefore says that the number of indices with \(|\mu_i|<c\) is
even.  On the other hand, Weyl's singular-value perturbation inequality
and \eqref{eq:x-c-second-singular} give, for the second-smallest singular
value \(s_{\rm next}\),
\[
 s_{\rm next}(x_*I_{2L-1}-A_{2L-1})
 \geq s_{L,2}^{(0)}-x_*>c,
\]
so at most one index can satisfy \(|\mu_i|<c\).  The number is thus
zero, and nonvanishing of the product excludes equality.  Hence
\(s_{\min}(x_*I_{2L-1}-A_{2L-1})>c_L\).  Since
\(0<x_*<2r\sin\frac{\pi}{2L}\), the point lies strictly in the positive
central gap, and \(d_{L-1}>c_L\).  This completes the proof of
\eqref{eq:central-interlace}.

\medskip\noindent\emph{Assembly of the global barrier inequality.}\quad
We finish the index comparison.  Put \(M=N+1\), so the larger and
smaller matrices have orders \(M-1\) and \(N-1\), respectively.  If
\(\Sigma=\diag(1,-1,1,-1,\ldots)\), then
\(\Sigma A_n\Sigma=-A_n\), so reflection reduces every noncentral
comparison to a positive gap.  By the interiority observation following the
definition of the gap heights, choose
\(x_j\in G_{M-1,j}^{\circ}\) at which its height is attained.  The hyperbolic
or elliptic construction above produces a point
\(\widehat x_j\in G_{N-1,j}\), and chord exhaustion gives
\[
 \gamma_{N-1,j}\geq g_{N-1}(\widehat x_j)
 >g_{M-1}(x_j)=\gamma_{M-1,j}.
\]
The reflected inequality handles the corresponding negative gap.

If \(M=2\ell_*+1\), the remaining single central height of the larger,
even-order matrix is \(c_{\ell_*}\), whereas each central height of the
smaller odd-order matrix is \(d_{\ell_*-1}>c_{\ell_*}\).  If
\(M=2\ell_*+2\), the two central heights of the larger odd-order matrix
equal \(d_{\ell_*}<c_{\ell_*}\), where \(c_{\ell_*}\) is the single
central height of the smaller even-order matrix.
Thus every height of the larger problem is strictly smaller than at
least one height of the smaller problem.  Choose \(j_*\) with
\(\gamma_{M-1,j_*}=\max_j\gamma_{M-1,j}\).  Then
\[
 \gamma_{M-1}=\gamma_{M-1,j_*}
 <\max_k\gamma_{N-1,k}=\gamma_{N-1}.
\]
Since \(M-1=n+1\) and \(N-1=n\), this proves
\(\gamma_{n+1}<\gamma_n\).
\end{proof}

\section{Topology, quantitative bounds, and the transition scale}
\label{sec:topology-bounds}

We begin with the standard component property for arbitrary matrices
\cite{BurkeLewisOverton2003}.  In the present setting it reduces the planar
topology to the real spectral interval.

\begin{lemma}\label{lem:component}
Every connected component of \(\sigma_\varepsilon(A_n)\) contains an
eigenvalue of \(A_n\).
\end{lemma}

\begin{proof}
The estimate
\[
 s_{\min}(zI_n-A_n)\geq |z|-\lVert A_n\rVert_2
\]
shows that the pseudospectrum is bounded.  Let \(\mathcal C\) be one of
its components.  If an eigenvalue \(\lambda\) belonged to
\(\overline{\mathcal C}\), continuity of the least singular value would
give a ball about \(\lambda\) contained in the pseudospectrum.  That
ball meets \(\mathcal C\), so maximality of connected components places
\(\lambda\) in \(\mathcal C\).

Suppose that \(\mathcal C\) contains no eigenvalue.  Its compact closure
is then disjoint from the spectrum.  On a neighbourhood of
\(\overline{\mathcal C}\), the function
\(\mathfrak u(z)=\lVert(zI_n-A_n)^{-1}\rVert_2\) is continuous and
subharmonic; see
\cite{TrefethenEmbree2005}.  Moreover,
\(\mathfrak u\leq\varepsilon^{-1}\) on \(\partial\mathcal C\), since a strict
inequality would put a neighbourhood of that boundary point in the same
component.  The maximum principle gives
\(\mathfrak u\leq\varepsilon^{-1}\) throughout \(\mathcal C\), contrary to the
definition of the pseudospectrum.
\end{proof}

The diagonal gauge introduced in the introduction gives
\[
 A_n=\mathcal G_n(r)\bigl[r(V_n+V_n^T)\bigr]
 \mathcal G_n(r)^{-1},
\]
and hence
\begin{equation}\label{eq:spectrum}
 \sigma(A_n)=\left\{2r\cos\frac{k\pi}{n+1}:1\leq k\leq n\right\}.
\end{equation}
Put
\[
 \Omega=\sigma_\varepsilon(A_n),\qquad
 \Omega_\R=\Omega\cap\R=\{x\in\R:g_n(x)<\varepsilon\}.
\]
Because \(A_n\) is real, the least singular value is invariant under
complex conjugation.  \Cref{lem:vertical} therefore implies that
\[
 \mathcal H_t(x+\ii y)=x+\ii(1-t)y,\qquad0\leq t\leq1,
\]
maps \(\Omega\) into itself.  If \(\mathcal C\) is a component of
\(\Omega\), the path \(t\mapsto\mathcal H_t(z)\) starts in
\(\mathcal C\) and hence remains there.  Thus \(\mathcal H_t\) retracts
\(\mathcal C\) onto
\(\mathcal C\cap\R\).

By \cref{lem:component}, this real intersection is nonempty.  It is
connected, being the continuous image \(\mathcal H_1(\mathcal C)\), and hence
lies in one interval component \(\mathcal K\) of \(\Omega_\R\).
Conversely, \(\mathcal K\) is a
connected subset of \(\Omega\) meeting \(\mathcal C\), so
\(\mathcal K\subset\mathcal C\).  Therefore
\(\mathcal C\cap\R=\mathcal K\), and the
components of \(\Omega\) correspond bijectively to the interval
components of \(\Omega_\R\).  In particular, each \(\mathcal C\)
retracts onto
an interval and is contractible.

If \(\Omega\) is connected, its real intersection is an interval
containing every eigenvalue, and therefore contains their convex hull
\(I_n^{\mathrm{sp}}\).  Conversely, if
\(I_n^{\mathrm{sp}}\subset\Omega\), that interval belongs to one
component and contains every eigenvalue; \cref{lem:component} rules out
any other component.  Compactness of \(I_n^{\mathrm{sp}}\) now gives
\begin{equation}\label{eq:connected-iff-gamma}
 \Omega\text{ is connected}
 \quad\Longleftrightarrow\quad
 I_n^{\mathrm{sp}}\subset\Omega
 \quad\Longleftrightarrow\quad
 \max_{x\in I_n^{\mathrm{sp}}}g_n(x)<\varepsilon
 \quad\Longleftrightarrow\quad
 \gamma_n<\varepsilon.
\end{equation}
The strict inequality reflects the open definition of the
pseudospectrum; equality leaves at least one real barrier point outside.

We now give quantitative estimates.  Put \(N=n+1\) and recall
\(\mathcal L_n(a)\) and \(\mathcal U_n(a)\) from
\eqref{eq:main-L-def}--\eqref{eq:main-U-def}.

\begin{proposition}[Sharp-order gap bounds]\label{prop:gap-bounds}
For every \(n\geq2\),
\begin{equation}\label{eq:gamma-LU}
 \mathcal L_n(a)\leq\gamma_n\leq\mathcal U_n(a).
\end{equation}
In particular, both bounds have order \((n+1)a^{n/2}\), with constants
depending only on \(a\).
\end{proposition}

\begin{proof}
For \(x=2r\cos\theta\in I_n^{\mathrm{sp}}\), set
\[
 \boldsymbol v=(v_j)_{j=1}^n,\qquad
 v_j=r^{j-1}\sin(j\theta),\qquad1\leq j\leq n.
\]
All entries of \((xI_n-A_n)\boldsymbol v\), except the last, vanish, and
the last one
is \(r^n\sin(N\theta)\).  Thus
\[
 g_n(2r\cos\theta)
 \leq\frac{r^n|\sin(N\theta)|}
 {\left(\sum_{j=1}^nr^{2j-2}\sin^2(j\theta)\right)^{1/2}}
 \leq r^n\frac{|\sin(N\theta)|}{\sin\theta}.
\]
Since \(\pi/N\leq\theta\leq(N-1)\pi/N\), this proves the upper bound.

For the lower bound, put
\(p_\theta=re^{\ii\theta}\), \(q_\theta=re^{-\ii\theta}\), so that
\(p_\theta+q_\theta=x\) and \(p_\theta q_\theta=a\).  For
\(\boldsymbol u=(u_j)_{j=1}^n\in\C^n\), put \(u_0=u_{n+1}=0\), and
define maps from \(\C^n\) to \(\C^{n+1}\) by
\[
 (\mathsf{C}_\theta\boldsymbol u)_j=p_\theta u_j-u_{j+1},\qquad
 (\mathsf{D}_\theta\boldsymbol u)_j=u_j-\overline{q_\theta}u_{j+1},
 \qquad0\leq j\leq n.
\]
For \(1\leq k\leq n\), direct multiplication gives
\[
 (\mathsf{D}_\theta^*\mathsf{C}_\theta\boldsymbol u)_k
 =(\mathsf{C}_\theta\boldsymbol u)_k
 -q_\theta(\mathsf{C}_\theta\boldsymbol u)_{k-1}
 =(p_\theta+q_\theta)u_k-u_{k+1}-p_\theta q_\theta u_{k-1}.
\]
Hence
\begin{equation}\label{eq:rectangular-factorization}
 xI_n-A_n=\mathsf{D}_\theta^*\mathsf{C}_\theta.
\end{equation}
The two Gram matrices are
diagonally unitarily similar to the tridiagonal Toeplitz matrix with
diagonal \(1+r^2\) and off-diagonal \(-r\).  Its eigenvalues are
\(1+r^2-2r\cos\frac{k\pi}{N}\), \(1\leq k\leq n\).  Consequently, for
every \(\boldsymbol u\in\C^n\),
\begin{equation}\label{eq:rectangular-min}
 \begin{aligned}
  \lVert\mathsf C_\theta\boldsymbol u\rVert_2
  &\geq \left(1+r^2-2r\cos\frac{\pi}{N}\right)^{1/2}
    \lVert\boldsymbol u\rVert_2,\\
  \lVert\mathsf D_\theta\boldsymbol u\rVert_2
  &\geq \left(1+r^2-2r\cos\frac{\pi}{N}\right)^{1/2}
    \lVert\boldsymbol u\rVert_2.
 \end{aligned}
\end{equation}
Put
\[
 \mathcal S_N=\sum_{j=0}^nr^{2j}=\frac{1-r^{2N}}{1-r^2}.
\]
Unit normal vectors to the two hyperplane ranges are
\[
 \boldsymbol u_C=\mathcal S_N^{-1/2}
 (\overline{p_\theta}^{\,n},\overline{p_\theta}^{\,n-1},\ldots,1)^T,
 \qquad
 \boldsymbol u_D=\mathcal S_N^{-1/2}(1,q_\theta,\ldots,q_\theta^n)^T.
\]
In the polar decompositions
\[
 \mathsf{C}_\theta=Q_C(\mathsf{C}_\theta^*\mathsf{C}_\theta)^{1/2},
 \qquad
 \mathsf{D}_\theta=Q_D(\mathsf{D}_\theta^*\mathsf{D}_\theta)^{1/2},
\]
the matrices \(Q_C,Q_D\) are isometries onto these ranges.  Two
hyperplanes in \(\C^{n+1}\) intersect in dimension \(n-1\), so
the cross-Gram matrix \(Q_D^*Q_C\) has \(n-1\) singular values equal to
\(1\); the remaining one is the cosine of the remaining principal
angle, \(|\langle\boldsymbol u_D,\boldsymbol u_C\rangle|\).  The
geometric sum gives
\begin{align}
 |\langle\boldsymbol u_D,\boldsymbol u_C\rangle|
 &=\frac{r^n}{\mathcal S_N}
 \left|\sum_{j=0}^{N-1}e^{\ii(2j-(N-1))\theta}\right|\notag\\
 &=
 \frac{r^n(1-r^2)}{1-r^{2N}}
 \left|\frac{\sin(N\theta)}{\sin\theta}\right|
 =:\chi_N(\theta).
\label{eq:chi}
\end{align}
Finally,
\[
 xI_n-A_n=(\mathsf{D}_\theta^*\mathsf{D}_\theta)^{1/2}
 Q_D^*Q_C(\mathsf{C}_\theta^*\mathsf{C}_\theta)^{1/2}.
\]
Using \(s_{\min}(XYZ)\geq
s_{\min}(X)s_{\min}(Y)s_{\min}(Z)\) and
\eqref{eq:rectangular-min} gives
\[
 g_n(2r\cos\theta)\geq
 \left(1+r^2-2r\cos\frac{\pi}{N}\right)\chi_N(\theta).
\]
Putting \(\theta=3\pi/(2N)\), which lies in the first spectral gap,
proves the lower bound.

Finally,
\[
 \mathcal U_n(a)\leq\frac12(n+1)r^n,
\]
while
\begin{equation}\label{eq:L-simple}
 \mathcal L_n(a)\geq C_{\rm lb}(a)(n+1)r^n,\qquad
 C_{\rm lb}(a)=\frac{2(1-r)^2(1-r^2)}{3\pi}.
\end{equation}
These follow from
\(\sin\frac{\pi}{N}\geq2/N\),
\(1+r^2-2r\cos\frac{\pi}{N}\geq(1-r)^2\), and
\(\sin\frac{3\pi}{2N}\leq3\pi/(2N)\).  This proves the asserted sharp
order.
\end{proof}

\begin{proof}[Proof of \cref{thm:main}]
By \cref{prop:gap-bounds},
\(\gamma_n\leq\mathcal U_n(a)\to0\).  \Cref{lem:mesh} says that
\(\gamma_n\) is strictly decreasing.  Equation
\eqref{eq:connected-iff-gamma} therefore shows that, once the
pseudospectrum is connected, it stays connected in every larger
dimension.  This proves \(N_e=N_f\).  Contractibility of every component
was proved immediately before \eqref{eq:connected-iff-gamma}.

If \(\mathcal L_n(a)\geq\varepsilon\), then
\(\gamma_n\geq\varepsilon\), so the pseudospectrum in dimension \(n\)
is disconnected.  If \(\mathcal U_n(a)<\varepsilon\), then
\(\gamma_n<\varepsilon\), so it is connected.  Strict decrease of
\(\gamma_n\) now gives \eqref{eq:N-two-sided-main}.

The comparison-envelope threshold, rather than \(N_c\) itself, can be
inverted in closed form.
Put
\[
 \Theta_n=(n+1)r^n,\qquad
 \mathcal T_r(t)=\min\{N\geq2:\Theta_m<t\text{ for every }m\geq N\}.
\]
Equations \eqref{eq:L-simple} and
\(\mathcal U_n(a)\leq\Theta_n/2\) give
\begin{equation}\label{eq:tail-threshold-bounds}
 \boxed{
 \mathcal T_r(\varepsilon/C_{\rm lb}(a))\leq N_c(a,\varepsilon)
 \leq\mathcal T_r(2\varepsilon).}
\end{equation}
Indeed, \eqref{eq:connected-iff-gamma} gives
\[
 N_c(a,\varepsilon)=
 \min\{N\geq2:\gamma_m<\varepsilon\text{ for every }m\geq N\}.
\]
The lower estimate \(C_{\rm lb}(a)\Theta_m\leq\gamma_m\) forces
\(\Theta_m<\varepsilon/C_{\rm lb}(a)\) throughout the tail beginning
at \(N_c\).  Conversely, \(\Theta_m<2\varepsilon\) for every
\(m\geq\mathcal T_r(2\varepsilon)\), together with
\(\gamma_m\leq\Theta_m/2\), forces connectedness throughout that tail.
Thus the polynomial factor \(n+1\), as well as the exponential factor
\(a^{n/2}\), is sharp on both sides.

\medskip\noindent\emph{Lambert-\(W\) inversion of the envelope threshold.}\quad
Let
\[
 \beta=-\log r=\frac12\log(1/a).
\]
For \(0<t<3r^2=\Theta_2\), let \(x_t>2\) be the solution on the decreasing
branch of \((x+1)e^{-\beta x}=t\).  This branch exists because
\[
 0<\beta rt<3\beta r^3\leq e^{-1};
\]
the last inequality is \(ue^{-u}\leq e^{-1}\) with \(u=3\beta\).
Moreover, \(t<\Theta_2\) places the descending root strictly beyond
\(x=2\).
Putting \(y_t=x_t+1\), the equation becomes
\((-\beta y_t)e^{-\beta y_t}=-\beta rt\), whence
\(y_t=-\beta^{-1}W_{-1}(-\beta rt)\).  Therefore
\begin{equation}\label{eq:tail-threshold-Lambert}
 \mathcal T_r(t)=\left\lfloor
 -\frac1\beta W_{-1}(-\beta rt)\right\rfloor,
\end{equation}
where \(W_{-1}\) is the lower real branch of the Lambert \(W\)-function.
Indeed, the continuous function \((x+1)r^x\) is unimodal and, after its
descending root, is strictly below \(t\).  Hence the integer tail
condition is \(N>x_t\); its least solution is
\(\lfloor x_t\rfloor+1=\lfloor x_t+1\rfloor\), including when \(x_t\)
is an integer.

The standard lower-branch expansion \cite{CorlessEtAl1996}
\[
 -W_{-1}(-u)=\log(1/u)+\log\!\log(1/u)+O(1),
 \qquad u\downarrow0,
\]
combined with \eqref{eq:tail-threshold-bounds} proves, for each fixed
\(a\in(0,1)\), as \(\varepsilon\downarrow0\),
\begin{equation}\label{eq:N-asymptotic}
 \boxed{
 N_c(a,\varepsilon)=
 \frac{2}{\log(1/a)}
 \left(\log(1/\varepsilon)+
       \log\!\log(1/\varepsilon)\right)+O_a(1).}
\end{equation}
The remainder is not asserted to be uniform as \(a\uparrow1\).

\medskip\noindent\emph{Exact two-site endpoint.}\quad
For $|x|\leq r$,
\[
 \det\bigl((xI_2-A_2)^*(xI_2-A_2)-a^2I_2\bigr)
 =x^2(x^2-2a-2a^2)\leq0.
\]
The matrix inside the determinant is Hermitian, and the nonpositive
product of its two eigenvalues forces its least eigenvalue to be
nonpositive.  Thus $s_{\min}(xI_2-A_2)\leq a$.  At $x=0$, the singular
values of $-A_2$ are $1$ and $a$, so equality is attained and
$\gamma_2=a$.  Hence, by \eqref{eq:connected-iff-gamma},
\begin{equation}\label{eq:N-two}
 N_c(a,\varepsilon)=2\quad\Longleftrightarrow\quad\varepsilon>a.
\end{equation}
This proves all assertions.
\end{proof}

\section{Numerical illustrations and reproducibility}\label{sec:numerics}

The computations in this section illustrate the theorem but are not used in
its proof.  The complete script
\texttt{make\_hatano\_nelson\_figures.py}, available in the public repository
cited in the Declarations, uses NumPy, SciPy, and Matplotlib.  It constructs
\(A_n(a)\) directly, computes Cartesian pseudospectral contours by singular
value decomposition, and obtains gapwise estimates of \(\gamma_n\) by bounded
scalar optimization.  As a separate check, let \(M_{\Delta x}\) be the largest
exact least singular value sampled on uniform real-gap grids of maximum
spacing \(\Delta x\).  The 1-Lipschitz property of
\(x\mapsto s_{\min}(xI_n-A_n)\) gives the rigorous implication
\[
 M_{\Delta x}\leq\gamma_n\leq M_{\Delta x}+\frac{\Delta x}{2}
\]
for exact samples.  The repository implementation applies this test to
NumPy/SciPy float64 SVD values, but it does not enclose SVD or eigenvalue
roundoff.  Its topology labels are therefore reproducible floating-point
checks, conditional on the computed samples, rather than validated interval
certificates.  They avoid reliance on pixel connectivity or the scalar
optimizer and are unused in every proof.

\begin{figure}[t]
 \centering
 \includegraphics[width=\textwidth]{fig1_connectedness_transition}
 \caption{Computed pseudospectral transition for \(a=1/4\) and
 \(\varepsilon=10^{-2}\).  Shading denotes the strict pseudospectrum, and
 crosses mark open-boundary eigenvalues.  The float64 barriers give the
 reported threshold \(N_c=7\); these checks are illustrative rather than
 interval certificates.}
 \label{fig:transition}
\end{figure}

The gapwise estimates are
\[
\begin{aligned}
 (\widehat\gamma_5,\widehat\gamma_6)
 &=(0.02909836,0.01465829),\\
 (\widehat\gamma_7,\widehat\gamma_8,\widehat\gamma_9)
 &=(0.007364219,0.003820033,0.001992747).
\end{aligned}
\]
The grid/Lipschitz check puts the computed values for \(n=5,6\) above
\(\varepsilon\) and those for \(n=7,8,9\) below it.  Thus the first two panels
are mode-resolved, and the remaining three form a single cloud.  This is a
clean-chain realization of the \(\widetilde N\) crossover described by
Kiorpelidis and Makris \cite{KiorpelidisMakris2025}.  Whenever the true
pseudospectrum is connected, every vertical section is an interval and the
real trace contains the full open-boundary spectral segment.

\begin{figure}[t]
 \centering
 \includegraphics[width=.76\textwidth]{fig2_gap_barrier_bounds}
 \caption{Gapwise float64 barrier estimates and the analytic enclosure
 \(\mathcal L_n(a)\leq\gamma_n\leq\mathcal U_n(a)\) for \(a=1/4\).
 The estimates first fall below \(\varepsilon=10^{-2}\) at \(n=7\); the
 exact bounds give \(6\leq N_c\leq9\), and all curves share the scale
 \((n+1)a^{n/2}\).}
 \label{fig:barriers}
\end{figure}

\Cref{fig:barriers} separates the two roles of the analysis.  The
strict comparison theorem identifies a unique crossing of a fixed
\(\varepsilon\), while the elementary upper and lower envelopes prove
its sharp exponential-polynomial scale without numerically resolving
the maximizing gap.

\begin{figure}[t]
 \centering
 \includegraphics[width=\textwidth]{fig3_hatano_nelson_physics}
 \caption{Real-axis barrier mechanism and skin-mode envelope for
 \(a=1/4\) and \(\varepsilon=10^{-2}\).  Panels (a)--(b) show \(g_n(x)\),
 the level \(\varepsilon\), and the real pseudospectrum for \(n=6,7\);
 panel (c) shows right eigenmodes for \(n=20\), with amplitude decay length
 \(\ell_{\rm skin}=2/\log(1/a)\).}
 \label{fig:physics}
\end{figure}

\Cref{fig:physics} places the topological criterion and its nonnormal
mechanism side by side.  The left panels show the final real-axis barrier
crossing the fixed uncertainty level; the right panel shows
the imaginary-gauge envelope responsible for the exponentially small
barrier scale.  Reversing the site convention reverses the localization edge
but not the connectedness criterion.

\clearpage
\section{Hatano--Nelson interpretation: a spectral-resolution site count}
\label{sec:physical-conclusion}

To connect the dimensionless matrix to the usual second-quantized
notation, consider
\[
 \widehat H_{\rm OBC}=-\sum_{j=1}^{n-1}
 \left(t_L\widehat c_j^\dagger\widehat c_{j+1}
 +t_R\widehat c_{j+1}^\dagger\widehat c_j\right),
 \qquad t_L>t_R>0.
\]
Its one-particle matrix is \(-t_LA_n(a)\), \(a=t_R/t_L\).  The
conventional minus sign is immaterial here because
\(\Sigma A_n(a)\Sigma=-A_n(a)\) for
\(\Sigma=\diag(1,-1,\ldots,(-1)^{n-1})\).  Put
\[
 t_0=\sqrt{t_Lt_R},\qquad
 h=\frac12\log(t_L/t_R)
 =\frac12\log(1/a).
\]
For the unitarily equivalent positive-sign representative
\(H_n=t_LA_n(a)\), multiplying the dimensionless gauge relation by \(t_L\)
gives the eigenvalues, unnormalized right eigenvectors \(\psi_m^R\) of
\(H_n\), and unnormalized column eigenvectors \(\psi_m^L\) of
\(H_n^\dagger\):
\begin{align*}
 E_m&=2t_0\cos\frac{m\pi}{n+1},\\
 \psi_m^R(j)&=e^{-h(j-1)}\sin\frac{jm\pi}{n+1},\\
 \psi_m^L(j)&=e^{h(j-1)}\sin\frac{jm\pi}{n+1},
 \qquad 1\leq j,m\leq n.
\end{align*}
With this normalization,
\((\psi_m^L)^\dagger\psi_{m'}^R=(n+1)\delta_{mm'}/2\).
These mode profiles are closely related to the original Hatano--Nelson
eigenfunction analysis \cite{HatanoNelson1998}.
Thus the amplitude skin depth of the right eigenmodes is
\[
 \ell_{\rm skin}=h^{-1}=\frac{2}{\log(1/a)},
\]
the amplitude e-folding length.  The intensity or probability envelope has
e-folding length \((2h)^{-1}=\ell_{\rm skin}/2\), while the exponential gauge
factors cancel in the corresponding biorthogonal left--right product.  The
gauge condition number is
\(\kappa_2(\mathcal G_n(e^{-h}))
=e^{h(n-1)}=a^{-(n-1)/2}\).

\subsection{Why a finite threshold exists}

The least singular value has the backward-error interpretation
\[
 s_{\min}(zI_n-A_n)
 =\min\{\lVert\Delta A\rVert_2:
 z\in\sigma(A_n+\Delta A)\}.
\]
Consequently, \(\gamma_n\) is not merely an auxiliary quantity from the
proof.  It is the largest, over the real interval between the extreme
open-boundary eigenvalues, of the minimum perturbation required to place an
eigenvalue at the specified point.  It is therefore the final barrier that
the uncertainty level must exceed to make the pseudospectrum connected.

For the dimensional positive-sign Hamiltonian \(H_n=t_LA_n(a)\), define
the dimensional barrier
\[
 \varepsilon_*(n)=t_L\gamma_n.
\]
Then the exact physical form of the connectedness criterion is
\[
 \sigma_{\varepsilon_{\rm phys}}(H_n)\text{ is connected}
 \quad\Longleftrightarrow\quad
 \varepsilon_{\rm phys}>\varepsilon_*(n).
\]
The strict decrease \(\gamma_{n+1}<\gamma_n\), together with
\(\gamma_n\to0\), explains why there is one finite threshold: a fixed
uncertainty level exceeds the decreasing barrier sequence once, after which
the cloud remains connected.  Equivalently, if
\(E_{\min}\) and \(E_{\max}\) are the extreme eigenvalues of \(H_n\), then
\[
 \varepsilon_*(n)^{-1}
 =\min_{x\in[E_{\min},E_{\max}]}
 \lVert(xI_n-H_n)^{-1}\rVert_2,
\]
with the resolvent norm interpreted as infinite at an eigenvalue.  At each
\(z\), this norm is the pointwise worst-case frequency-domain amplification
from forcing to response; for \(H_n\) it has physical units of inverse
energy.  Thus \(\varepsilon_*(n)^{-1}\) is the minimum pointwise amplification
over the spectral interval, while \(\gamma_n^{-1}\) is the corresponding
dimensionless minimum for \(A_n\).  At and beyond \(N_c\), this minimum
exceeds \(\varepsilon_{\rm phys}^{-1}\).

The Hatano--Nelson content lies in the rate at which this barrier collapses.
From \cref{prop:gap-bounds},
\[
 \varepsilon_*(n)
 \asymp_a t_L(n+1)e^{-hn}
 =t_0(n+1)e^{-h(n-1)}
 =t_0(n+1)e^{-(n-1)/\ell_{\rm skin}}.
\]
The nonunitary imaginary gauge becomes exponentially ill-conditioned across
the chain, and right and left skin modes accumulate at opposite edges.
Accordingly,
\[
 \gamma_n^{-1}\asymp_a\frac{e^{hn}}{n+1},
\]
so gauge conditioning and worst-case spectral sensitivity have the same
exponential rate \(h\), up to fixed-\(a\) constants and the displayed
polynomial factor.  Skin localization and barrier collapse are parallel
consequences of the same imaginary gauge.

Let \(\delta=\varepsilon_{\rm phys}/t_L\).  Balancing the fixed uncertainty
level with the barrier scale above gives, as \(\delta\downarrow0\),
\[
 N_c(a,\delta)=\ell_{\rm skin}
\left(\log(1/\delta)+\log\!\log(1/\delta)\right)+O_a(1).
\]
Here \(N_c\) is a number of sites.  If the lattice spacing is \(d\), the
corresponding end-to-end chain length is approximately \((N_c-1)d\); only in
lattice-spacing units is it convenient to call \(N_c\) a
spectral-resolution length.  Measured in amplitude skin depths, its leading
value is
\(N_c/\ell_{\rm skin}=\log(1/\delta)+\log\!\log(1/\delta)+O_a(1)\): this is
the site count needed for a fixed uncertainty level to exceed the final
barrier.  Improving the resolution by a factor \(q\) increases the leading
resolvable site count by \(\ell_{\rm skin}\log q\).  At fixed
\(\delta=\varepsilon_{\rm phys}/t_L\), stronger nonreciprocity shortens
\(\ell_{\rm skin}\) and hence lowers the leading small-\(\delta\) threshold.
This is an asymptotic statement, not a finite-parameter monotonicity theorem
in \(a\).

The exact two-site endpoint makes the same interpretation especially
transparent.  Since \(\gamma_2=a\),
\[
 \varepsilon_*(2)=t_La=t_R,
 \qquad
 N_c(a,\varepsilon_{\rm phys}/t_L)=2
 \quad\Longleftrightarrow\quad
 \varepsilon_{\rm phys}>t_R.
\]
Indeed, for the positive-sign dimer,
\[
 H_2=\begin{bmatrix}0&t_L\\t_R&0\end{bmatrix},
 \qquad
 \Delta H=-t_R\mathbf e_2\mathbf e_1^T,
 \qquad
 H_2+\Delta H=\begin{bmatrix}0&t_L\\0&0\end{bmatrix}.
\]
The norm-\(t_R\) perturbation cancels the weaker directed hopping and
produces a defective double eigenvalue at zero.  Hence, if the uncertainty
level exceeds \(t_R\), mode resolution is lost already in the shortest
nontrivial chain; no long-chain buildup is needed.  At exact equality the
needed perturbation is excluded by the strict convention
\(\lVert\Delta H\rVert_2<\varepsilon_{\rm phys}\): the component closures
touch at zero but the strict pseudospectrum is disconnected (indeed,
\(N_c(a,a)=3\)).
Here loss of mode resolution means loss of component-based resolution under
unstructured worst-case spectral-norm uncertainty.  The displayed rank-one
hopping perturbation is a special structured realization; the theorem does
not assert that generic onsite or hopping disorder realizes it.

Kiorpelidis and Makris \cite{KiorpelidisMakris2025} denoted by
\(\widetilde N\) the observed fixed-\(\varepsilon\) crossover from separate
eigenvalue-centered clouds to a single cloud, and associated it with the end
of exponential growth of the pseudospectral abscissa.  For the clean open
Hatano--Nelson chain, \(N_c\) is the rigorous clean-chain analogue of that
crossover.  Here \(N_c\) is defined by strict connectedness of the open
Euclidean pseudospectrum, whereas \(\widetilde N\) is their observed
cloud-merger crossover; the exact barrier criterion and strict barrier
decrease additionally rule out re-entrant splitting.  Their convention
\(A=-\ii H\) rotates
the pseudospectrum by ninety degrees and does not change connectedness.

\subsection{What is lost at the threshold}

The relevant uncertainty is deterministic worst-case Hamiltonian or model
uncertainty, not a quantum uncertainty principle.  A point
\(z\in\sigma_\varepsilon(A_n)\) records only that \emph{some} perturbation of
norm below \(\varepsilon\) can place an eigenvalue at \(z\); the point does
not come tagged with that perturbation or with a parent eigenvalue.
Connected components nevertheless give a precise, perturbation-stable
coarse label.  Define two unperturbed eigenvalues to be in the same
\(\varepsilon\)-resolution class when they lie in the same component of
\(\sigma_\varepsilon(A_n)\).  For any continuous perturbation path
\(\Delta A(s)\) satisfying \(\lVert\Delta A(s)\rVert_2<\varepsilon\), every
continuous eigenvalue branch of \(A_n+\Delta A(s)\) remains inside
\(\sigma_\varepsilon(A_n)\) and therefore cannot pass between two distinct
components.  Equivalently, each component carries a constant total algebraic
multiplicity of perturbed eigenvalues along such a path.

When a component contains one unperturbed eigenvalue, this class gives an
individual-mode label.  Partial mergers can combine some labels before the
full threshold, so \(N_c\) need not be the first size at which one pair of
eigenvalues becomes ambiguous.  Rather, \(N_c\) is the \emph{last} merger:
for \(n<N_c\) at least two robust spectral classes remain, whereas for
\(n\geq N_c\) the sole component contains all \(n\) eigenvalues.  The
location alone then no longer determines a perturbation-independent,
topologically protected open-boundary spectral cluster.  If the actual
perturbation or its continuation path is known, eigenvalue branches may
still be tracked.

This statement concerns a union over perturbations.  It does not assert that
one perturbation places eigenvalues at every point of a connecting region.
Nor does it describe probabilities of observing different locations.  In
resolvent language, for \(z\notin\sigma(A)\),
\[
 \lVert(zI_n-A)^{-1}\rVert_2
 =\max_{\boldsymbol f\neq0}
 \frac{\lVert\boldsymbol u\rVert_2}{\lVert\boldsymbol f\rVert_2},
 \qquad (zI_n-A)\boldsymbol u=\boldsymbol f,
\]
so the single cloud is a region of large pointwise resolvent amplification at the specified
worst-case resolution, rather than the spectrum of one noisy sample.

\subsection{Scope and conclusion}

The finite-size existence of a merger is not exclusive to the skin effect.
At the normal endpoint \(a=1\), increasingly dense Hermitian eigenvalues also
merge at fixed \(\varepsilon\), but on the algebraic level-spacing scale.
What is specific to this Hatano--Nelson family is the precise barrier scale
\((n+1)a^{n/2}\), up to fixed-\(a\) constants, its exact connectedness
consequence and resulting logarithmic plus log-log threshold law, and the
appearance of \(\ell_{\rm skin}\) as the leading coefficient.  This also
explains why the limit \(a\uparrow1\) in the physical threshold asymptotic
above is not uniform.

Restricting perturbations to onsite, hopping, boundary, or Toeplitz
structure can produce a smaller attainable set and a different merger
threshold, and requires a structured-pseudospectral analysis.
Toeplitz-structured pseudospectra are treated in
\cite{BottcherGrudsky2005,ButtaGuglielmiNoschese2012}.  Robustness under
random imperfections is studied in \cite{AmmariEtAl2024Stability};
pseudospectral boundary sensitivity in correlated or disordered systems in
\cite{OkumaSato2021}; and complex and structured-pseudospectral scaling in
asymmetric lattices in \cite{KiorpelidisMakris2025}.  Likewise, the
single-cloud threshold is a spectral-resolution crossover, not the end of
dynamical exponential sensitivity; Kiorpelidis and Makris explicitly show
that long-time dynamical amplification can persist beyond their
\(\widetilde N\).

In summary, vertical monotonicity reduces planar connectedness to a
backward-error barrier on the real axis, and strict barrier decrease makes
its crossing with a fixed uncertainty level unique.  The imaginary gauge
sets the exponentially small scale of that barrier, while the skin depth
sets the leading finite-size resolution site count.  For the clean open
Hatano--Nelson chain, \(N_c\) therefore marks the global loss of robust
spectral-cluster separation under unstructured worst-case uncertainty,
without conflating that loss with the skin effect itself, a particular
structured perturbation, or dynamical amplification.

\section*{Declarations}

\textbf{Use of generative AI.}
Generative AI tools were used to assist with manuscript drafting and editing,
proof auditing, Lean formalization, writing and checking the reproducibility
code used to generate the numerical figures, and literature and citation
checking.  The author independently verified the mathematical statements,
computations, code, figures, and references.  The author assumes
responsibility for all content.

\textbf{Manuscript, code, and formalization availability.}
The source and compiled PDF of this manuscript, the Python figure-generation
script and its reproduction instructions, and the complete pinned Lean~4
formalization with its manuscript-to-formalization map and audit
documentation are available in the public GitHub repository
\url{https://github.com/jinshanmu/Hatano-Nelson}.

\phantomsection
\raggedbottom
\bibliographystyle{siamplain}
\bibliography{final_connected_pseudospectrum_proof}

\end{document}
